\documentclass[10pt,a4paper]{article}
\usepackage[T1]{fontenc}
\usepackage[utf8]{inputenc}
\usepackage{lmodern}
\usepackage[margin=22mm,headheight=14pt]{geometry}
\usepackage{microtype}
\usepackage{amsmath,amssymb}
\usepackage{longtable,booktabs,array,calc}
\usepackage{graphicx}
\usepackage{xcolor}
\usepackage{xurl}
\usepackage{fancyhdr}
\usepackage{enumitem}
\usepackage{needspace}
\usepackage{fvextra}
\usepackage{hyperref}
\definecolor{reportblue}{HTML}{164B70}
\hypersetup{colorlinks=true,linkcolor=reportblue,urlcolor=reportblue,
  pdftitle={Qwen3.5 Inference on the SpaceMiT
K1},pdfauthor={riscv-accl project}}
\urlstyle{tt}
\setlength{\parindent}{0pt}
\setlength{\parskip}{5pt}
\setlength{\emergencystretch}{3em}
\setlength{\LTpre}{8pt}
\setlength{\LTpost}{8pt}
\setcounter{secnumdepth}{-1}
\setcounter{tocdepth}{1}
\renewcommand{\arraystretch}{1.15}
\setlist{nosep,leftmargin=1.5em}
\pagestyle{fancy}
\fancyhf{}
\fancyhead[L]{\small\textcolor{reportblue}{riscv-accl / SpaceMiT K1}}
\fancyhead[R]{\small 2026-10-12}
\fancyfoot[C]{\small\thepage}
\providecommand{\tightlist}{\setlength{\itemsep}{0pt}\setlength{\parskip}{0pt}}
\usepackage{color}
\usepackage{fancyvrb}
\newcommand{\VerbBar}{|}
\newcommand{\VERB}{\Verb[commandchars=\\\{\}]}
\DefineVerbatimEnvironment{Highlighting}{Verbatim}{commandchars=\\\{\}}
% Add ',fontsize=\small' for more characters per line
\newenvironment{Shaded}{}{}
\newcommand{\AlertTok}[1]{\textcolor[rgb]{1.00,0.00,0.00}{\textbf{#1}}}
\newcommand{\AnnotationTok}[1]{\textcolor[rgb]{0.38,0.63,0.69}{\textbf{\textit{#1}}}}
\newcommand{\AttributeTok}[1]{\textcolor[rgb]{0.49,0.56,0.16}{#1}}
\newcommand{\BaseNTok}[1]{\textcolor[rgb]{0.25,0.63,0.44}{#1}}
\newcommand{\BuiltInTok}[1]{#1}
\newcommand{\CharTok}[1]{\textcolor[rgb]{0.25,0.44,0.63}{#1}}
\newcommand{\CommentTok}[1]{\textcolor[rgb]{0.38,0.63,0.69}{\textit{#1}}}
\newcommand{\CommentVarTok}[1]{\textcolor[rgb]{0.38,0.63,0.69}{\textbf{\textit{#1}}}}
\newcommand{\ConstantTok}[1]{\textcolor[rgb]{0.53,0.00,0.00}{#1}}
\newcommand{\ControlFlowTok}[1]{\textcolor[rgb]{0.00,0.44,0.13}{\textbf{#1}}}
\newcommand{\DataTypeTok}[1]{\textcolor[rgb]{0.56,0.13,0.00}{#1}}
\newcommand{\DecValTok}[1]{\textcolor[rgb]{0.25,0.63,0.44}{#1}}
\newcommand{\DocumentationTok}[1]{\textcolor[rgb]{0.73,0.13,0.13}{\textit{#1}}}
\newcommand{\ErrorTok}[1]{\textcolor[rgb]{1.00,0.00,0.00}{\textbf{#1}}}
\newcommand{\ExtensionTok}[1]{#1}
\newcommand{\FloatTok}[1]{\textcolor[rgb]{0.25,0.63,0.44}{#1}}
\newcommand{\FunctionTok}[1]{\textcolor[rgb]{0.02,0.16,0.49}{#1}}
\newcommand{\ImportTok}[1]{#1}
\newcommand{\InformationTok}[1]{\textcolor[rgb]{0.38,0.63,0.69}{\textbf{\textit{#1}}}}
\newcommand{\KeywordTok}[1]{\textcolor[rgb]{0.00,0.44,0.13}{\textbf{#1}}}
\newcommand{\NormalTok}[1]{#1}
\newcommand{\OperatorTok}[1]{\textcolor[rgb]{0.40,0.40,0.40}{#1}}
\newcommand{\OtherTok}[1]{\textcolor[rgb]{0.00,0.44,0.13}{#1}}
\newcommand{\PreprocessorTok}[1]{\textcolor[rgb]{0.74,0.48,0.00}{#1}}
\newcommand{\RegionMarkerTok}[1]{#1}
\newcommand{\SpecialCharTok}[1]{\textcolor[rgb]{0.25,0.44,0.63}{#1}}
\newcommand{\SpecialStringTok}[1]{\textcolor[rgb]{0.73,0.40,0.53}{#1}}
\newcommand{\StringTok}[1]{\textcolor[rgb]{0.25,0.44,0.63}{#1}}
\newcommand{\VariableTok}[1]{\textcolor[rgb]{0.10,0.09,0.49}{#1}}
\newcommand{\VerbatimStringTok}[1]{\textcolor[rgb]{0.25,0.44,0.63}{#1}}
\newcommand{\WarningTok}[1]{\textcolor[rgb]{0.38,0.63,0.69}{\textbf{\textit{#1}}}}
\RecustomVerbatimEnvironment{Highlighting}{Verbatim}{breaklines=true,fontsize=\small,commandchars=\\\{\}}
\begin{document}
\begin{titlepage}
\vspace*{15mm}
{\large\textcolor{reportblue}{RISC-V LOCAL INFERENCE}}\par
\vspace{12mm}
{\Huge\bfseries Qwen3.5 Inference on the SpaceMiT K1\par}
\vspace{8mm}
{\Large Infrastructure Optimization and Performance Evaluation\par}
\vspace{12mm}
{\large Evidence cutoff: October 12, 2026 (Asia/Singapore)\par}
\vspace{14mm}
\begin{tabular}{@{}ll@{}}
\textbf{Platform} & MUSE-Pi-Pro, K1/X60, 16 GiB RAM \\
\textbf{Models} & Qwen3.5 2B and 4B \\
\textbf{Implementation} & SpaceMiT llama.cpp fork \\
\textbf{Package} & Technical report and evidence appendix \\
\end{tabular}
\vfill
\textcolor{reportblue}{\rule{\textwidth}{1pt}}\par
This report separates verified release results, later candidates, and prediction limits.
Its measurements do not establish a cumulative project speedup.\par
\vspace{6mm}
\href{https://github.com/MoyaMryia/riscv-accl}{github.com/MoyaMryia/riscv-accl}
\end{titlepage}
{\small\tableofcontents}
\clearpage
\textbf{Project report - evidence cutoff: October 12, 2026,
Asia/Singapore}\\
Platform: MUSE-Pi-Pro, SpaceMiT K1/X60, 16 GiB RAM\\
Models: Qwen3.5 2B and 4B conversational checkpoints\\
Implementation: SpaceMiT's open-source llama.cpp fork

\hypertarget{abstract}{%
\section{Abstract}\label{abstract}}

This project implements and tests local Qwen3.5 inference on the
MUSE-Pi-Pro with the SpaceMiT K1/X60 and 16 GiB RAM. The main release
retains Q4\_0 weights, F16 K/V, complete inputs, and direct decoding.
The 256-bit RVV attention path reduces 8,192-token first-token latency
by 49.2\% on 2B and 64.0\% on 4B. Separate IME1 routing tests reduce
64-token decode time by 14.40\% and 28.62\%. These comparisons use
different controls and do not establish a cumulative speedup.

Later candidates improve selected prefill or decode workloads on the
optimized baseline. Their complete-answer tests contain 48 identical
pairs, but shared model errors remain. The three-candidate combination
fails one decode gate with a 3.0220477\% regression against the
attention-only candidate. Its second round and complete-answer stages do
not run. The verified October 7 release remains the pinned package;
routing stays opt-in.

The project also provides a portable performance evaluator and native
operator data. V7 covers all nine reserved groups with wide intervals,
but its point directions agree with measurements in only 9/27 cases. The
first operator batch contains 1,344 observations across 32
configurations. These training data do not establish independent
prediction accuracy. K1 remote access fails before a confirmed second
batch. Recovery remains pending; this document does not authorize new
board work.

\hypertarget{report-status-and-evidence-levels}{%
\section{Report status and evidence
levels}\label{report-status-and-evidence-levels}}

\begin{longtable}[]{@{}ll@{}}
\toprule
\begin{minipage}[b]{0.47\columnwidth}\raggedright
Evidence level\strut
\end{minipage} & \begin{minipage}[b]{0.47\columnwidth}\raggedright
Meaning in this report\strut
\end{minipage}\tabularnewline
\midrule
\endhead
\begin{minipage}[t]{0.47\columnwidth}\raggedright
Implemented\strut
\end{minipage} & \begin{minipage}[t]{0.47\columnwidth}\raggedright
Source code or an executable tool exists. This status does not establish
a performance gain.\strut
\end{minipage}\tabularnewline
\begin{minipage}[t]{0.47\columnwidth}\raggedright
Verified\strut
\end{minipage} & \begin{minipage}[t]{0.47\columnwidth}\raggedright
The stated checks pass for the named configuration and records.\strut
\end{minipage}\tabularnewline
\begin{minipage}[t]{0.47\columnwidth}\raggedright
Failed gate\strut
\end{minipage} & \begin{minipage}[t]{0.47\columnwidth}\raggedright
A measured result does not meet its declared criterion. Later versions
do not remove that result.\strut
\end{minipage}\tabularnewline
\begin{minipage}[t]{0.47\columnwidth}\raggedright
Pending\strut
\end{minipage} & \begin{minipage}[t]{0.47\columnwidth}\raggedright
A measurement or independent check is incomplete.\strut
\end{minipage}\tabularnewline
\begin{minipage}[t]{0.47\columnwidth}\raggedright
Not implemented\strut
\end{minipage} & \begin{minipage}[t]{0.47\columnwidth}\raggedright
The project does not provide the stated capability.\strut
\end{minipage}\tabularnewline
\bottomrule
\end{longtable}

This report separates the verified release, later candidates, and
performance-model work. Dated reports describe their original campaigns.
The
\href{https://github.com/MoyaMryia/riscv-accl/blob/main/spacemit/reports/SUBMISSION-EVIDENCE.md}{evidence
appendix} maps claims to records and tools. The report and appendix form
one review package. No external competition submission is recorded.

\hypertarget{objective-and-scope}{%
\section{1. Objective and scope}\label{objective-and-scope}}

The task deploys both required Qwen3.5 models locally through the
SpaceMiT fork. It targets prefill, decode, and long-context inference.
Deliverables include baseline comparisons, executable code, and
reproducible test results.

The current optimization mission concerns infrastructure. It retains
complete inputs, Q4\_0 weights and F16 K/V. Retrieval, prompt
shortening, token pruning and model replacement are excluded. Earlier
Q8\_0 and speculative experiments are included as supporting evidence
with their separate settings and limits. All model generation runs
locally. An optional cloud judge is used only for evaluation of
completed answers, including cache reuse and MTP comparisons; inference
does not require that judge.

\hypertarget{requirement-coverage}{%
\subsection{Requirement coverage}\label{requirement-coverage}}

\begin{longtable}[]{@{}ll@{}}
\toprule
\begin{minipage}[b]{0.47\columnwidth}\raggedright
Requirement\strut
\end{minipage} & \begin{minipage}[b]{0.47\columnwidth}\raggedright
Evidence and boundary\strut
\end{minipage}\tabularnewline
\midrule
\endhead
\begin{minipage}[t]{0.47\columnwidth}\raggedright
Both 2B and 4B models\strut
\end{minipage} & \begin{minipage}[t]{0.47\columnwidth}\raggedright
Both have local inference, matched prefill comparisons and routing
decode comparisons.\strut
\end{minipage}\tabularnewline
\begin{minipage}[t]{0.47\columnwidth}\raggedright
SpaceMiT open-source llama.cpp\strut
\end{minipage} & \begin{minipage}[t]{0.47\columnwidth}\raggedright
Current release pins official-fork \nolinkurl{a990751} and measured
source hashes. The historical patch stack targets
\nolinkurl{5ad05d8}.\strut
\end{minipage}\tabularnewline
\begin{minipage}[t]{0.47\columnwidth}\raggedright
Baseline and optimized results\strut
\end{minipage} & \begin{minipage}[t]{0.47\columnwidth}\raggedright
Matched RVV off/on and GEMM staging/bypass tables below; each comparison
states its baseline.\strut
\end{minipage}\tabularnewline
\begin{minipage}[t]{0.47\columnwidth}\raggedright
Quantifiable infrastructure improvement\strut
\end{minipage} & \begin{minipage}[t]{0.47\columnwidth}\raggedright
RVV reduces long-prompt TTFT; the new routing screen improves both
phases on both models.\strut
\end{minipage}\tabularnewline
\begin{minipage}[t]{0.47\columnwidth}\raggedright
Deployment, scripts and tests\strut
\end{minipage} & \begin{minipage}[t]{0.47\columnwidth}\raggedright
Integration guide, patch helper, native harnesses, tmux runners and
archived evidence are provided.\strut
\end{minipage}\tabularnewline
\begin{minipage}[t]{0.47\columnwidth}\raggedright
Generation above 1 token/s\strut
\end{minipage} & \begin{minipage}[t]{0.47\columnwidth}\raggedright
Demonstrated in short/moderate-context decode tests; this is not
guaranteed at every long context.\strut
\end{minipage}\tabularnewline
\begin{minipage}[t]{0.47\columnwidth}\raggedright
Cross-architecture comparison\strut
\end{minipage} & \begin{minipage}[t]{0.47\columnwidth}\raggedright
Excluded from the release requirements at the user's direction.\strut
\end{minipage}\tabularnewline
\bottomrule
\end{longtable}

The
\href{https://github.com/MoyaMryia/riscv-accl/blob/main/spacemit/reports/2026-09-27-next-steps-and-measures.md}{original
requested checklist} and
\href{https://github.com/MoyaMryia/riscv-accl/blob/main/env/competition-background.md}{competition
background} have different scopes. This report states the measured
coverage rather than treating every suggested advanced challenge as
completed.

\hypertarget{platform-models-and-measurement-method}{%
\section{2. Platform, models and measurement
method}\label{platform-models-and-measurement-method}}

The board runs Bianbu 2.3.5 with a RISC-V Linux environment. The
evaluated backend uses RVV and SpaceMiT IME1; four inference workers are
used because the IME-capable cores are 0--3. Build instructions specify
GCC 14 and the matching SPERT runtime. The latest runtime audit records
the performance governor at 1.6 GHz. Runtime libraries and original
build flags are preserved in experiment provenance.

The principal benchmark files are:

\begin{itemize}
\tightlist
\item
  \nolinkurl{Qwen3.5-2B-MTP-Q4_0-embQ4_0-dv64k.gguf}
\item
  \nolinkurl{Qwen3.5-4B-MTP-Q4_0-embQ4_0-dv64k.gguf}
\end{itemize}

These are the project's conversational checkpoint derivatives; their
filenames omit the word \nolinkurl{Instruct}. MTP tensors are present,
but the principal RVV and routing results use \textbf{direct decoding}.
Weight and cache precision are distinct: Q4\_0 weights with F16 K/V is
the common current configuration.

\hypertarget{controls-and-metrics}{%
\subsection{Controls and metrics}\label{controls-and-metrics}}

\begin{longtable}[]{@{}ll@{}}
\toprule
\begin{minipage}[b]{0.47\columnwidth}\raggedright
Item\strut
\end{minipage} & \begin{minipage}[b]{0.47\columnwidth}\raggedright
Protocol\strut
\end{minipage}\tabularnewline
\midrule
\endhead
\begin{minipage}[t]{0.47\columnwidth}\raggedright
Model/runtime controls\strut
\end{minipage} & \begin{minipage}[t]{0.47\columnwidth}\raggedright
Same model and full prompt within each comparison; four workers,
batch/microbatch 32, flash attention; direct greedy decoding.\strut
\end{minipage}\tabularnewline
\begin{minipage}[t]{0.47\columnwidth}\raggedright
Run order\strut
\end{minipage} & \begin{minipage}[t]{0.47\columnwidth}\raggedright
Four cold ABBA requests for each routing comparison: control, candidate,
candidate, control. Earlier campaign repetitions are stated
separately.\strut
\end{minipage}\tabularnewline
\begin{minipage}[t]{0.47\columnwidth}\raggedright
Cache isolation\strut
\end{minipage} & \begin{minipage}[t]{0.47\columnwidth}\raggedright
Cold requests require zero reused prompt tokens and the complete
intended input count.\strut
\end{minipage}\tabularnewline
\begin{minipage}[t]{0.47\columnwidth}\raggedright
TTFT\strut
\end{minipage} & \begin{minipage}[t]{0.47\columnwidth}\raggedright
Client time to first token; includes work visible to the requesting
client.\strut
\end{minipage}\tabularnewline
\begin{minipage}[t]{0.47\columnwidth}\raggedright
Prefill time\strut
\end{minipage} & \begin{minipage}[t]{0.47\columnwidth}\raggedright
Server prompt-processing time; a one-token request isolates this phase
but does not establish answer quality.\strut
\end{minipage}\tabularnewline
\begin{minipage}[t]{0.47\columnwidth}\raggedright
Decode rate/time\strut
\end{minipage} & \begin{minipage}[t]{0.47\columnwidth}\raggedright
Server generation measurements after prefill; latest screen captures all
64 requested tokens.\strut
\end{minipage}\tabularnewline
\begin{minipage}[t]{0.47\columnwidth}\raggedright
Memory\strut
\end{minipage} & \begin{minipage}[t]{0.47\columnwidth}\raggedright
Whole-process RSS is distinguished from estimated KV payload.\strut
\end{minipage}\tabularnewline
\begin{minipage}[t]{0.47\columnwidth}\raggedright
Correctness\strut
\end{minipage} & \begin{minipage}[t]{0.47\columnwidth}\raggedright
Native finite-output/guard checks plus input, token and text hashes,
completion counts and actual kernel activation logs.\strut
\end{minipage}\tabularnewline
\bottomrule
\end{longtable}

For the routing screen, a candidate's mean time reduction must exceed
both 3\% and the observed full control range. Required comparisons must
have no clear regression. This is a predeclared engineering gate, not a
statistical significance test. Two samples per arm do not establish a
population confidence interval. Exact output agreement is limited to the
tested inputs and settings.

Time reduction and throughput increase use different denominators. For
control time \(t_b\) and candidate time \(t_c\), the time reduction is

\[R_t = 100\left(1-\frac{t_c}{t_b}\right)\%. \]

For equal output counts, the throughput increase is

\[G_q = 100\left(\frac{t_b}{t_c}-1\right)\%. \]

The tables state whether they use means, medians, or median paired
ratios. A ratio of medians need not equal the median paired ratio.

\hypertarget{implemented-optimizations}{%
\section{3. Implemented optimizations}\label{implemented-optimizations}}

\hypertarget{rvv-attention-for-the-k1s-vector-width}{%
\subsection{3.1 RVV attention for the K1's vector
width}\label{rvv-attention-for-the-k1s-vector-width}}

The original wide-vector prototype did not activate on the K1. The
subsequent 256-bit implementation enables the F16-KV wide-head prefill
path and is packaged as optional patch 0009. It is enabled with
\nolinkurl{SPINE_FA_WIDE_TILE=1}; activation is checked in the server
logs. It improves prompt processing and does not establish a
single-token decode gain. Quantized K/V and other model head shapes are
outside its demonstrated dispatch.

\hypertarget{production-q4_0-gemm-routing}{%
\subsection{3.2 Production Q4\_0 GEMM
routing}\label{production-q4_0-gemm-routing}}

Production attribution identified substantial staging-copy work: 9.55\%
/ 9.93\% of summed instrumented prefill worker elapsed time for 2B / 4B.
These sums include overlapping waits and instrumentation overhead, so
they are not removable wall-time estimates. The new code instead tests
the existing direct-memory branch in production
\nolinkurl{forward_mul_mat}.

Eligibility is restricted to Q4\_0 K32/N16 traits using IME1, with K
divisible by 32 and N divisible by 16. The selected buffer is set to
null with size zero; weight packing, activation quantization and IME
arithmetic are retained. The generated library is isolated from the
original build.

\begin{longtable}[]{@{}lll@{}}
\toprule
\begin{minipage}[b]{0.30\columnwidth}\raggedright
\nolinkurl{SPINE_K1_GEMM_ROUTE}\strut
\end{minipage} & \begin{minipage}[b]{0.30\columnwidth}\raggedright
Selection\strut
\end{minipage} & \begin{minipage}[b]{0.30\columnwidth}\raggedright
Validation status\strut
\end{minipage}\tabularnewline
\midrule
\endhead
\begin{minipage}[t]{0.30\columnwidth}\raggedright
Unset, invalid or 0\strut
\end{minipage} & \begin{minipage}[t]{0.30\columnwidth}\raggedright
Original staging\strut
\end{minipage} & \begin{minipage}[t]{0.30\columnwidth}\raggedright
Control\strut
\end{minipage}\tabularnewline
\begin{minipage}[t]{0.30\columnwidth}\raggedright
1\strut
\end{minipage} & \begin{minipage}[t]{0.30\columnwidth}\raggedright
Bypass for M \textgreater{} 1\strut
\end{minipage} & \begin{minipage}[t]{0.30\columnwidth}\raggedright
Numerical, operator and model prefill gates passed\strut
\end{minipage}\tabularnewline
\begin{minipage}[t]{0.30\columnwidth}\raggedright
2\strut
\end{minipage} & \begin{minipage}[t]{0.30\columnwidth}\raggedright
Bypass for M = 1\strut
\end{minipage} & \begin{minipage}[t]{0.30\columnwidth}\raggedright
Numerical, operator and 64-token model decode gates passed\strut
\end{minipage}\tabularnewline
\begin{minipage}[t]{0.30\columnwidth}\raggedright
3\strut
\end{minipage} & \begin{minipage}[t]{0.30\columnwidth}\raggedright
Bypass for both\strut
\end{minipage} & \begin{minipage}[t]{0.30\columnwidth}\raggedright
Combined 8k phase gates pass on both models; overall quality gate
remains open\strut
\end{minipage}\tabularnewline
\bottomrule
\end{longtable}

Mode 2 can affect single-row work during prefill as well as decoding.
The current TCM API reports unavailable/fake/zero geometry, while SPERT
supplies non-null 128 KiB worker buffers. Their physical placement is
unverified; the measured result concerns staging policy under this
runtime.

\hypertarget{earlier-kernel-and-runtime-work}{%
\subsection{3.3 Earlier kernel and runtime
work}\label{earlier-kernel-and-runtime-work}}

The packaged work also includes recurrent-state in-place write-back,
microbatch selection, Q8\_0 IME1 support, Q4\_0 M4 scale construction
and workload-specific speculative policies. Q8\_0 IME1 raised 4B pp64
from 1.117 to 7.472 token/s and tg32 from 0.866 to 1.349 token/s in its
separate higher-precision comparison. Q4\_0 remained faster overall.

M4 scale construction raised 4B Q4\_0 pp128 from 9.058 to 9.360 token/s
(3.34\%). Corrected analysis gives a conditional independent-launch
Welch 95\% interval of +0.095 to +0.511 token/s, with only two launches
per arm. The roughly 1.8\% English/Chinese mapped-MTP gains have
intervals including zero. These small effects remain workload-specific.
Detailed implementation, attribution and statistics are linked in the
evidence appendix.

\hypertarget{matched-performance-results}{%
\section{4. Matched performance
results}\label{matched-performance-results}}

\hypertarget{rvv-attention-versus-the-generic-attention-path}{%
\subsection{4.1 RVV attention versus the generic attention
path}\label{rvv-attention-versus-the-generic-attention-path}}

Both arms use the same model and configuration within each comparison,
Q4\_0 weights, F16 K/V, direct decoding and 32 generated tokens. The
baseline here has RVV wide attention disabled.

\begin{longtable}[]{@{}lrrrr@{}}
\toprule
Model / input tokens & Generic TTFT & RVV TTFT & TTFT reduction &
Launches per arm\tabularnewline
\midrule
\endhead
2B / 2,048 & 113.78 s & 92.45 s & 18.7\% & 2\tabularnewline
4B / 2,048 & 343.95 s & 231.54 s & 32.7\% & 2\tabularnewline
2B / 8,192 & 862.62 s & 438.05 s & 49.2\% & 2\tabularnewline
4B / 8,192 & 3,194.99 s & 1,150.42 s & 64.0\% & 2\tabularnewline
2B / 16,384 & 2,808.28 s & 1,124.57 s & 60.0\% & 1\tabularnewline
4B / 16,384 & 11,062.98 s & 3,216.23 s & 70.9\% & 1\tabularnewline
\bottomrule
\end{longtable}

All these comparisons have complete matching token/text hashes. The
2k/8k rows use the isolated build; 16k uses the integrated build. The
single-pair 16k observations need replication. Decode and peak RSS were
effectively unchanged within these attention comparisons.

A single RVV-enabled 2B request with an actual 32,768-token prompt
completed at 3,291.21 s TTFT, 9.96 prefill token/s and 2,768.4 MiB peak
RSS. Its 32 output tokens completed; decode was 0.95 token/s. This
establishes bounded feasibility, with no matched 32k control or
answer-quality claim. 4B/32k and 64k are unmeasured.

\hypertarget{new-gemm-routing-versus-the-current-optimized-baseline}{%
\subsection{4.2 New GEMM routing versus the current optimized
baseline}\label{new-gemm-routing-versus-the-current-optimized-baseline}}

The baseline already includes accepted wide attention and layout 0. Each
row has two launches per arm. Prefill rows use mode 1 and one output
token; decode rows use mode 2, a full 256-token prompt and 64 captured
output tokens.

\begin{longtable}[]{@{}lrrr@{}}
\toprule
Model / measured phase & Control mean & Candidate mean & Time
reduction\tabularnewline
\midrule
\endhead
2B / 512-token prefill & 22.052917 s & 20.625322 s &
6.47\%\tabularnewline
4B / 512-token prefill & 56.740043 s & 52.304415 s &
7.82\%\tabularnewline
2B / 2,048-token prefill & 92.634959 s & 86.709042 s &
\textbf{6.40\%}\tabularnewline
4B / 1,024-token prefill & 115.551047 s & 106.596561 s &
\textbf{7.75\%}\tabularnewline
2B / 64-token decode & 17.016656 s & 14.566911 s &
\textbf{14.40\%}\tabularnewline
4B / 64-token decode & 43.553046 s & 31.087887 s &
\textbf{28.62\%}\tabularnewline
\bottomrule
\end{longtable}

Decode throughput rose from \textbf{3.76 to 4.39 token/s} for 2B and
\textbf{1.47 to 2.06 token/s} for 4B. Mode 2 showed no clear prefill
regression. All six comparisons cleared their declared gates. These are
time reductions; throughput increases have a different percentage
denominator.

The board run completed successfully in 38.16 minutes. Validation
included 104 real production graph cases per arm across five arms, 144
production operator records, and 24 model requests. Every model
comparison matched input/output hashes and uncached counts. All 236
collected artifact hashes were independently verified, and end-of-run
checks matched 46 original source/object/build files. The verifier
recomputes the required gates from individual records.

These initial routing modes were tested separately. Their gains must not
be added or combined with older RVV percentages to claim cumulative
acceleration. The subsequent
\href{https://github.com/MoyaMryia/riscv-accl/blob/main/spacemit/reports/2026-10-03-staged-validation-results.md}{staged
validation} completed in 4.16 hours with 83 artifact hashes verified.
Combined mode 3 at 8192 input/64 output tokens reduces prefill/decode
time by 4.43\%/9.49\% on 2B and 5.79\%/15.06\% on 4B. All eight requests
match hashes and cold counts. The 12 full-document answers stop
naturally and match paired baseline text. Both 2B arms lack required
tracing citations, so the overall gate remains open. The protocol skips
4B/32k. These are bounded engineering pilots.

\hypertarget{independent-memory-bandwidth-and-routing-confirmation}{%
\subsection{4.3 Independent memory-bandwidth and routing
confirmation}\label{independent-memory-bandwidth-and-routing-confirmation}}

The
\href{https://github.com/MoyaMryia/riscv-accl/blob/main/spacemit/reports/2026-10-06-k1-roofline-results.md}{October
6 campaign} completed 150 native full-read scans, 24 cold
direct-inference requests and two separate prefill profiles. The
four-core model-sized streaming reference is approximately 7 GB/s. Main
weight-read proxies, excluding MTP draft tensors, are 1.060 GB for 2B
and 2.369 GB for 4B, giving weight-only references of 6.60/2.94 tok/s.
These are traffic references rather than measured DDR utilization or
attainable inference guarantees.

At 256 prompt tokens, route-0/route-3 decode medians are 3.79/4.48 tok/s
for 2B and 1.48/2.06 for 4B; median paired throughput gains are
18.01\%/38.70\%. At 2,048 tokens the medians are 3.31/3.83 and
1.22/1.62, with paired gains 15.47\%/32.36\%. All requests generate the
same 64 capped output tokens and match their paired control. This
confirms the existing routing mechanism; these gains must not be
multiplied by previous routing measurements.

The profiles identify convolution and synchronization as investigation
targets, but sampled prefill CPU shares are not wall-time savings or
decode attribution. This campaign adds no useful-answer qualification
and does not demonstrate that the hardware limit has been reached.

\hypertarget{later-candidates-and-the-stopped-combination}{%
\section{5. Later candidates and the stopped
combination}\label{later-candidates-and-the-stopped-combination}}

\hypertarget{independent-candidates-on-the-optimized-baseline}{%
\subsection{5.1 Independent candidates on the optimized
baseline}\label{independent-candidates-on-the-optimized-baseline}}

These comparisons retain the fixed 2B/4B GGUF files, four workers, batch
and microbatch 32, F16 K/V, route 3, temperature 0, and seed 42. MTP is
disabled. Prefill tests use 512 or 2,048 input tokens. The ranges below
describe observed configurations across two rounds. They are not
population confidence intervals.

\begin{longtable}[]{@{}llll@{}}
\toprule
\begin{minipage}[b]{0.22\columnwidth}\raggedright
Candidate\strut
\end{minipage} & \begin{minipage}[b]{0.22\columnwidth}\raggedright
Change\strut
\end{minipage} & \begin{minipage}[b]{0.22\columnwidth}\raggedright
Measured model result\strut
\end{minipage} & \begin{minipage}[b]{0.22\columnwidth}\raggedright
Scope\strut
\end{minipage}\tabularnewline
\midrule
\endhead
\begin{minipage}[t]{0.22\columnwidth}\raggedright
C1\strut
\end{minipage} & \begin{minipage}[t]{0.22\columnwidth}\raggedright
Hybrid convolution; original path for small batches\strut
\end{minipage} & \begin{minipage}[t]{0.22\columnwidth}\raggedright
One-output-token TTFT falls 4.26--6.53\%\strut
\end{minipage} & \begin{minipage}[t]{0.22\columnwidth}\raggedright
64 cold requests across two rounds; 24 complete answers\strut
\end{minipage}\tabularnewline
\begin{minipage}[t]{0.22\columnwidth}\raggedright
C2-r2\strut
\end{minipage} & \begin{minipage}[t]{0.22\columnwidth}\raggedright
GDN register kernel reduces intermediate state access\strut
\end{minipage} & \begin{minipage}[t]{0.22\columnwidth}\raggedright
One-output-token TTFT falls 4.35--5.95\%\strut
\end{minipage} & \begin{minipage}[t]{0.22\columnwidth}\raggedright
96 cold requests across the screen and two rounds; 24 complete
answers\strut
\end{minipage}\tabularnewline
\begin{minipage}[t]{0.22\columnwidth}\raggedright
C3\strut
\end{minipage} & \begin{minipage}[t]{0.22\columnwidth}\raggedright
C1 with the older C2\strut
\end{minipage} & \begin{minipage}[t]{0.22\columnwidth}\raggedright
TTFT falls 8.35--9.73\% against its disabled control; 3.53--4.63\%
against C1\strut
\end{minipage} & \begin{minipage}[t]{0.22\columnwidth}\raggedright
128 cold requests across two controls and two rounds; 24 complete
answers\strut
\end{minipage}\tabularnewline
\begin{minipage}[t]{0.22\columnwidth}\raggedright
B1-r2\strut
\end{minipage} & \begin{minipage}[t]{0.22\columnwidth}\raggedright
Shared attention path with a fallback for small work groups\strut
\end{minipage} & \begin{minipage}[t]{0.22\columnwidth}\raggedright
At 2,048 input/65 output tokens, decode time falls 10.48--10.64\% on 2B
and 16.95--20.32\% on 4B\strut
\end{minipage} & \begin{minipage}[t]{0.22\columnwidth}\raggedright
64 cold requests across two rounds; 24 complete answers\strut
\end{minipage}\tabularnewline
\bottomrule
\end{longtable}

GDN means gated delta network, the recurrent-state component in these
models. C3 contains the older C2, not C2-r2. All four versions complete
96 answers with 48 identical text pairs. Shared arithmetic, citation,
code, and Chinese-language errors remain. These tests do not show an
answer-quality improvement.

C2-r2 increases 4B/512/65 decode time by 2.9480\% in its first round.
This result is near the 3\% regression gate and is not a decode gain.
The old controller deletes some native dumps before the later review.
Saved logs cannot replace a fresh comparison of those dumps.
Complete-answer requests also lack full operating-system thread
snapshots. The
\href{https://github.com/MoyaMryia/riscv-accl/blob/main/spacemit/reports/2026-10-09-local-candidate-result-rejudgment.md}{candidate
review} records these limits.

\hypertarget{combined-r1-a-valid-performance-stop}{%
\subsection{5.2 Combined-r1: a valid performance
stop}\label{combined-r1-a-valid-performance-stop}}

The combination uses one library with independent switches. C0 disables
the three new candidates but retains wide RVV attention and route 3. B
enables B1-r2 only. P enables C1 and C2-r2. F enables B1-r2, C1, and
C2-r2. C0 is an optimized control, not the initial vendor source.

The run stops on October 11, 2026, at 18:28 Asia/Singapore with exit
code 0 and state \nolinkurl{performance-gate-stopped}. The
\href{https://github.com/MoyaMryia/riscv-accl/blob/main/spacemit/reports/2026-10-12-combined-stop-audit.md}{October
12 audit} recomputes all 48 comparisons from the first round's 96
requests. It confirms zero reused prompt tokens, output receipts, input
hashes, and saved thread snapshots. One required comparison fails: F
versus B for 4B/2,048 input/65 output tokens. Decode time rises
3.0220477\%, above the declared 3\% gate. The full control range is
0.0726418\%.

Each table value is the mean of two requests in that arm. All times are
milliseconds.

\begin{longtable}[]{@{}lrrrr@{}}
\toprule
Model / input tokens & B prefill & F prefill & B decode & F
decode\tabularnewline
\midrule
\endhead
2B / 512 & 20959.891 & 18769.821 & 14913.766 & 14920.559\tabularnewline
2B / 2048 & 88041.229 & 79032.342 & 15716.909 & 15714.604\tabularnewline
4B / 512 & 52863.175 & 47537.975 & 30820.230 & 31336.530\tabularnewline
4B / 2048 & 226049.395 & 203459.167 & 33608.751 &
34624.424\tabularnewline
\bottomrule
\end{longtable}

For the failed configuration, F prefill time falls 9.9935\% against B. F
decode time also falls 17.5476\% against C0. Neither result clears the
required F-versus-B regression gate. The second round and
complete-answer tests do not run. The combination does not complete
acceptance.

The audit does not recollect or compare all 536 formal state dumps. It
checks saved input and text hashes, not reconstructed inputs or new
answers. Startup thread snapshots allow CPU 0--7 and do not establish
worker positions during decode. This gap does not prove incorrect core
placement. Both failed-comparison arms emit 2,048 shared-attention log
lines per request. The evidence does not establish extra attention logs
as the cause. Logging cost remains part of the measured result.

\hypertarget{local-revision-and-comparison-boundary}{%
\subsection{5.3 Local revision and comparison
boundary}\label{local-revision-and-comparison-boundary}}

The local dispatch revision changes one GDN condition order. It returns
single-token and non-128-state calls before it reads the prefill switch.
All 90 dispatch combinations select the same paths. The cross-compiler
syntax check and four local tool checks pass. Full numerical checks,
target assembly, operator tests, and model tests remain pending. This
revision is not a verified fix for the regression. The stopped
combination remains a failed result.

The
\href{https://github.com/MoyaMryia/riscv-accl/blob/main/spacemit/reports/2026-10-11-original-current-comparison.md}{original/current
comparison} contains historical phase measurements and explicit version
identities. Its October 11 pending entries predate the stop audit. Use
the October 12 audit for the combination's current status. No matched
test compares the initial vendor source directly with a final accepted
combination. The separate RVV, routing, and candidate gains must not be
added or multiplied.

\hypertarget{negative-results-and-correctness-limits}{%
\section{6. Negative results and correctness
limits}\label{negative-results-and-correctness-limits}}

\begin{longtable}[]{@{}lll@{}}
\toprule
\begin{minipage}[b]{0.30\columnwidth}\raggedright
Candidate\strut
\end{minipage} & \begin{minipage}[b]{0.30\columnwidth}\raggedright
Measured outcome\strut
\end{minipage} & \begin{minipage}[b]{0.30\columnwidth}\raggedright
Decision\strut
\end{minipage}\tabularnewline
\midrule
\endhead
\begin{minipage}[t]{0.30\columnwidth}\raggedright
Compact attention layouts\strut
\end{minipage} & \begin{minipage}[t]{0.30\columnwidth}\raggedright
Original 2k pilot gains of 0.89\% / 1.98\%; later fast-screen model
gains below 3\%\strut
\end{minipage} & \begin{minipage}[t]{0.30\columnwidth}\raggedright
Retain layout 0\strut
\end{minipage}\tabularnewline
\begin{minipage}[t]{0.30\columnwidth}\raggedright
Further IME scheduling / scale gathering\strut
\end{minipage} & \begin{minipage}[t]{0.30\columnwidth}\raggedright
Numerical checks pass; operator screens do not qualify\strut
\end{minipage} & \begin{minipage}[t]{0.30\columnwidth}\raggedright
Retain original schedule\strut
\end{minipage}\tabularnewline
\begin{minipage}[t]{0.30\columnwidth}\raggedright
Hybrid SSM convolution/RVV\strut
\end{minipage} & \begin{minipage}[t]{0.30\columnwidth}\raggedright
Exact graph/state checks and 24 natural answers pass; total latency
falls 2.92\% / 2.13\%, below the adoption gate\strut
\end{minipage} & \begin{minipage}[t]{0.30\columnwidth}\raggedright
Disabled in the release\strut
\end{minipage}\tabularnewline
\begin{minipage}[t]{0.30\columnwidth}\raggedright
Shared dense FFN activation packing\strut
\end{minipage} & \begin{minipage}[t]{0.30\columnwidth}\raggedright
Duplicate work is only 0.024--0.035\% of summed GEMM worker elapsed in
the decode audit\strut
\end{minipage} & \begin{minipage}[t]{0.30\columnwidth}\raggedright
Implementation declined\strut
\end{minipage}\tabularnewline
\begin{minipage}[t]{0.30\columnwidth}\raggedright
Fixed K32 M1 inner-loop specialization\strut
\end{minipage} & \begin{minipage}[t]{0.30\columnwidth}\raggedright
Exact raw/production/full-output checks pass; 144 operator samples fail
advancement\strut
\end{minipage} & \begin{minipage}[t]{0.30\columnwidth}\raggedright
Disabled; model stages skipped\strut
\end{minipage}\tabularnewline
\begin{minipage}[t]{0.30\columnwidth}\raggedright
Recurrent prefill fusion\strut
\end{minipage} & \begin{minipage}[t]{0.30\columnwidth}\raggedright
Operator time falls 24--28\%; 2B/512 model time falls 2.38\%, below
gate\strut
\end{minipage} & \begin{minipage}[t]{0.30\columnwidth}\raggedright
Experimental; no advancement\strut
\end{minipage}\tabularnewline
\begin{minipage}[t]{0.30\columnwidth}\raggedright
Direct V attention access\strut
\end{minipage} & \begin{minipage}[t]{0.30\columnwidth}\raggedright
Operator gains do not become model gains: 2B/2k is 1.67\% slower; 4B/1k
effectively unchanged\strut
\end{minipage} & \begin{minipage}[t]{0.30\columnwidth}\raggedright
No adoption\strut
\end{minipage}\tabularnewline
\begin{minipage}[t]{0.30\columnwidth}\raggedright
Larger QK grouping\strut
\end{minipage} & \begin{minipage}[t]{0.30\columnwidth}\raggedright
Operator qualification insufficient\strut
\end{minipage} & \begin{minipage}[t]{0.30\columnwidth}\raggedright
No model advancement\strut
\end{minipage}\tabularnewline
\begin{minipage}[t]{0.30\columnwidth}\raggedright
Occupied-page gathering\strut
\end{minipage} & \begin{minipage}[t]{0.30\columnwidth}\raggedright
Throughput falls about 2--3\%; backing allocation retained; resumed
telemetry shows no compacted spans\strut
\end{minipage} & \begin{minipage}[t]{0.30\columnwidth}\raggedright
No speed or full-paged-scheduling claim\strut
\end{minipage}\tabularnewline
\begin{minipage}[t]{0.30\columnwidth}\raggedright
OpenCL GPU offload\strut
\end{minipage} & \begin{minipage}[t]{0.30\columnwidth}\raggedright
Zero layers transferred; backend rejects the PowerVR device\strut
\end{minipage} & \begin{minipage}[t]{0.30\columnwidth}\raggedright
No GPU acceleration claim\strut
\end{minipage}\tabularnewline
\begin{minipage}[t]{0.30\columnwidth}\raggedright
Vulkan route\strut
\end{minipage} & \begin{minipage}[t]{0.30\columnwidth}\raggedright
Dropped before measured offload\strut
\end{minipage} & \begin{minipage}[t]{0.30\columnwidth}\raggedright
No GPU acceleration claim\strut
\end{minipage}\tabularnewline
\begin{minipage}[t]{0.30\columnwidth}\raggedright
Speculative timing gate / hybrid IME-RVV dispatch\strut
\end{minipage} & \begin{minipage}[t]{0.30\columnwidth}\raggedright
Regress tested workloads\strut
\end{minipage} & \begin{minipage}[t]{0.30\columnwidth}\raggedright
Excluded\strut
\end{minipage}\tabularnewline
\bottomrule
\end{longtable}

Long-code MTP outputs differ from direct outputs beginning at generated
index 179 for 2B and 303 for 4B. Smaller draft lengths and RS rollback
did not fix this. Full-vocabulary target verification does not prove
numerical identity between different execution paths. Historical
matching 128-token samples therefore do not justify calling MTP
universally lossless. Direct decoding is used for the principal results
here.

The newer 512-output checkpoint/RS samples match direct outputs on both
models. Native RS rollback replay fails strict logit checks. Each model
has 18 mismatches, including argmax changes. Full-state restore and
same-shape controls pass. This newer evidence narrows the next diagnosis
and does not establish general MTP identity or supersede the older
long-output traces.

The separate
\href{https://github.com/MoyaMryia/riscv-accl/blob/main/spacemit/reports/2026-10-03-mtp-usefulness-results.md}{complete-answer
MTP pilot} completed 24 requests with RS disabled. Total paired task
latency falls 8.63\% on 2B and 15.22\% on 4B; both models pass the
predeclared relative fact/code/judge quality gate. The 4B Unicode code
differs from direct but passes all 106 functional checks. Absolute
usefulness still fails shared interval-code errors and 2B arithmetic,
while Chinese prose becomes slower. Checkpoint MTP remains optional.
These single-pair task results support workload-specific usefulness, not
universal losslessness or general adoption.

The later
\href{https://github.com/MoyaMryia/riscv-accl/blob/main/spacemit/reports/2026-10-05-adaptive-mtp-results.md}{adaptive
infrastructure screen} completed 54 timing requests and verified 47
artifacts. Median paired code throughput rises 50.74\% on 2B and 53.35\%
on 4B. The 2B code is capped. The 4B code is shorter under adaptive MTP
and fails reconstruction in every arm. Prose/QA wall time increases
5.77\%/4.76\% on 2B and 2.97\%/3.62\% on 4B despite disabled drafting.
The relative judge gate passes on 4B, but both models fail absolute
correctness and overall timing gates. This establishes bounded
throughput, not faster useful answers. Adaptive MTP remains unqualified
and optional.

Two complete 4,096-token requests per model and direct/MTP mode
demonstrated bounded generation stability, including a roughly 43-minute
4B direct request. These earlier traces are not a long-duration
validation of the new routing candidate, and they do not establish
indefinite service stability.

\hypertarget{full-input-cache-reuse-separate-supporting-evidence}{%
\section{7. Full-input cache reuse: separate supporting
evidence}\label{full-input-cache-reuse-separate-supporting-evidence}}

Same-process prefix reuse preserves the complete document and sends
changed questions through a local server slot. It avoids recomputing an
unchanged prefix and is separate from the cold-kernel comparisons.
Cached and cold answers need not be token-identical; usefulness is
checked on complete answers.

The original alternating quality protocol evaluated both models, 4k/8k
documents and six questions: 24 naturally completed cached/cold pairs.
An optional MiMo judge scored each pair in both answer orders,
supplementing fact, citation, completion, cache and slot checks.

\begin{longtable}[]{@{}lrrl@{}}
\toprule
\begin{minipage}[b]{0.19\columnwidth}\raggedright
Model / document\strut
\end{minipage} & \begin{minipage}[b]{0.25\columnwidth}\raggedleft
Cold / cached mean score\strut
\end{minipage} & \begin{minipage}[b]{0.25\columnwidth}\raggedleft
Median cold / cached TTFT\strut
\end{minipage} & \begin{minipage}[b]{0.19\columnwidth}\raggedright
Descriptive decision\strut
\end{minipage}\tabularnewline
\midrule
\endhead
\begin{minipage}[t]{0.19\columnwidth}\raggedright
2B / 4k\strut
\end{minipage} & \begin{minipage}[t]{0.25\columnwidth}\raggedleft
4.50 / 4.50 out of 5\strut
\end{minipage} & \begin{minipage}[t]{0.25\columnwidth}\raggedleft
211.11 / 3.90 s\strut
\end{minipage} & \begin{minipage}[t]{0.19\columnwidth}\raggedright
Review required\strut
\end{minipage}\tabularnewline
\begin{minipage}[t]{0.19\columnwidth}\raggedright
2B / 8k\strut
\end{minipage} & \begin{minipage}[t]{0.25\columnwidth}\raggedleft
4.83 / 4.83\strut
\end{minipage} & \begin{minipage}[t]{0.25\columnwidth}\raggedleft
474.31 / 5.01 s\strut
\end{minipage} & \begin{minipage}[t]{0.19\columnwidth}\raggedright
Pilot pass\strut
\end{minipage}\tabularnewline
\begin{minipage}[t]{0.19\columnwidth}\raggedright
4B / 4k\strut
\end{minipage} & \begin{minipage}[t]{0.25\columnwidth}\raggedleft
4.83 / 4.83\strut
\end{minipage} & \begin{minipage}[t]{0.25\columnwidth}\raggedleft
551.91 / 10.80 s\strut
\end{minipage} & \begin{minipage}[t]{0.19\columnwidth}\raggedright
Pilot pass\strut
\end{minipage}\tabularnewline
\begin{minipage}[t]{0.19\columnwidth}\raggedright
4B / 8k\strut
\end{minipage} & \begin{minipage}[t]{0.25\columnwidth}\raggedleft
4.83 / 5.00\strut
\end{minipage} & \begin{minipage}[t]{0.25\columnwidth}\raggedleft
1,271.37 / 14.65 s\strut
\end{minipage} & \begin{minipage}[t]{0.19\columnwidth}\raggedright
Pilot pass\strut
\end{minipage}\tabularnewline
\bottomrule
\end{longtable}

One 2B/4k question scores 2/5 in both arms. The matrix has no universal
quality pass. That question does not show a cache-specific regression.
One judge and six questions per configuration cannot establish general
quality equivalence. These results also do not validate the new
GEMM-routing candidate. Length-capped one-, 32- or 64-token speed
samples are not complete-answer evidence.

\hypertarget{performance-evaluator-and-native-operator-data}{%
\section{8. Performance evaluator and native operator
data}\label{performance-evaluator-and-native-operator-data}}

\hypertarget{purpose-and-accuracy-boundary}{%
\subsection{8.1 Purpose and accuracy
boundary}\label{purpose-and-accuracy-boundary}}

The evaluator combines actual instruction traces with measured component
costs. Its portable examples retain frozen models, error bounds, and
independent test results. QEMU supplies instruction and numerical
evidence; it does not simulate K1 cycle timing. A component-cost
estimate is not a verified full-model latency prediction. All
performance candidates remain available in this evaluator. Its accuracy
checks do not use candidate elimination gates.

The
\href{https://github.com/MoyaMryia/riscv-accl/blob/main/spacemit/reports/2026-10-12-k1-model-v7.md}{V7
audit} fixes 28 configuration groups before measurements. It uses 19
groups to calibrate errors and reserves nine groups for independent
tests. The groups include 1, 2, and 4 accumulator chains with 0--9
independent integer instructions. Each configuration has two arms and 8,
32, or 128 unrolled iterations. Both arms also change input-register
reuse; instruction order is not the only difference.

\begin{longtable}[]{@{}ll@{}}
\toprule
\begin{minipage}[b]{0.47\columnwidth}\raggedright
V7 check\strut
\end{minipage} & \begin{minipage}[b]{0.47\columnwidth}\raggedright
Recorded result\strut
\end{minipage}\tabularnewline
\midrule
\endhead
\begin{minipage}[t]{0.47\columnwidth}\raggedright
Numerical scope\strut
\end{minipage} & \begin{minipage}[t]{0.47\columnwidth}\raggedright
168 native complete dumps match QEMU for the same ELF; all-one signed
8-bit inputs\strut
\end{minipage}\tabularnewline
\begin{minipage}[t]{0.47\columnwidth}\raggedright
Actual traces\strut
\end{minipage} & \begin{minipage}[t]{0.47\columnwidth}\raggedright
336 complete traces; vendor QEMU 10.0.2 executes the IME ELF\strut
\end{minipage}\tabularnewline
\begin{minipage}[t]{0.47\columnwidth}\raggedright
Native timing\strut
\end{minipage} & \begin{minipage}[t]{0.47\columnwidth}\raggedright
CPU 0; 512/2048/8192 loops; seven repetitions; raw user-cycle event
0x22\strut
\end{minipage}\tabularnewline
\begin{minipage}[t]{0.47\columnwidth}\raggedright
Samples\strut
\end{minipage} & \begin{minipage}[t]{0.47\columnwidth}\raggedright
2,394 calibration samples; 1,134 reserved-test samples\strut
\end{minipage}\tabularnewline
\begin{minipage}[t]{0.47\columnwidth}\raggedright
Frozen ratio-error factor\strut
\end{minipage} & \begin{minipage}[t]{0.47\columnwidth}\raggedright
2.1216883936561333\strut
\end{minipage}\tabularnewline
\begin{minipage}[t]{0.47\columnwidth}\raggedright
Interval coverage\strut
\end{minipage} & \begin{minipage}[t]{0.47\columnwidth}\raggedright
9/9 reserved groups and 27/27 unrolled versions\strut
\end{minipage}\tabularnewline
\begin{minipage}[t]{0.47\columnwidth}\raggedright
Point-direction agreement\strut
\end{minipage} & \begin{minipage}[t]{0.47\columnwidth}\raggedright
9/27 versions\strut
\end{minipage}\tabularnewline
\begin{minipage}[t]{0.47\columnwidth}\raggedright
Absolute gain error\strut
\end{minipage} & \begin{minipage}[t]{0.47\columnwidth}\raggedright
Median 29.540 percentage points; maximum 50.990 percentage points\strut
\end{minipage}\tabularnewline
\begin{minipage}[t]{0.47\columnwidth}\raggedright
Portable replay\strut
\end{minipage} & \begin{minipage}[t]{0.47\columnwidth}\raggedright
Identical key predictions on the local host and K1; 62 tool checks
pass\strut
\end{minipage}\tabularnewline
\bottomrule
\end{longtable}

Five two-chain reserved groups contain 15 versions with an incorrect
predicted direction. The single-chain measurements are near zero and do
not establish a statistically supported change. The reserved set
contains no four-chain group. Wide interval coverage does not establish
reliable point predictions or an optimization priority order. Nine
observed groups cannot confirm a population coverage rate of 95\%. The
stated coverage requires exchangeable groups within the fixed binary and
data scope. The target is the next paired fixture measurement, not a
noise-free hardware mean.

The all-one numerical fixture does not establish correctness for
arbitrary data. Upstream QEMU 11.1.1 rejects the ELF with an illegal
instruction; the failure record remains. Vendor QEMU executes the same
ELF with VLEN 256. Frequency snapshots read 1,600,000 kHz but do not
establish a constant frequency during each measurement. V7 replay takes
approximately 0.44 s locally and 8.11 s on K1 after integrity checks.
Those times measure the host tool, not the target algorithm. V7 server
replay remains pending.

V5 retains its 7/9 coverage gap. V6's 9/9 aligned-loop group coverage
does not transfer to mixed IME loops. Shared-cache traffic, multithread
synchronization costs, and full IME operator time remain pending. A gem5
backend calibrated on K1 is not implemented. Independent prediction
confidence for arbitrary new algorithms is not established.

\hypertarget{first-native-operator-training-batch}{%
\subsection{8.2 First native operator training
batch}\label{first-native-operator-training-batch}}

The
\href{https://github.com/MoyaMryia/riscv-accl/blob/main/spacemit/reports/2026-10-12-k1-operator-data.md}{operator
report} records a successful first batch. The run covers October 12,
2026, 09:24:08--09:25:44 Asia/Singapore. It holds the shared board lock
and exits with code 0. The local host and project server verify
identical summaries and receipts.

\begin{longtable}[]{@{}ll@{}}
\toprule
\begin{minipage}[b]{0.47\columnwidth}\raggedright
Batch item\strut
\end{minipage} & \begin{minipage}[b]{0.47\columnwidth}\raggedright
Verified scope\strut
\end{minipage}\tabularnewline
\midrule
\endhead
\begin{minipage}[t]{0.47\columnwidth}\raggedright
Shape groups\strut
\end{minipage} & \begin{minipage}[t]{0.47\columnwidth}\raggedright
16 training groups; eight per operator family\strut
\end{minipage}\tabularnewline
\begin{minipage}[t]{0.47\columnwidth}\raggedright
Configurations\strut
\end{minipage} & \begin{minipage}[t]{0.47\columnwidth}\raggedright
32; one-thread and four-thread settings per shape\strut
\end{minipage}\tabularnewline
\begin{minipage}[t]{0.47\columnwidth}\raggedright
Numerical dumps\strut
\end{minipage} & \begin{minipage}[t]{0.47\columnwidth}\raggedright
64; paired versions match byte for byte\strut
\end{minipage}\tabularnewline
\begin{minipage}[t]{0.47\columnwidth}\raggedright
Timing observations\strut
\end{minipage} & \begin{minipage}[t]{0.47\columnwidth}\raggedright
1,344; all retained\strut
\end{minipage}\tabularnewline
\begin{minipage}[t]{0.47\columnwidth}\raggedright
IME worker observations\strut
\end{minipage} & \begin{minipage}[t]{0.47\columnwidth}\raggedright
1,680\strut
\end{minipage}\tabularnewline
\begin{minipage}[t]{0.47\columnwidth}\raggedright
File receipts\strut
\end{minipage} & \begin{minipage}[t]{0.47\columnwidth}\raggedright
327 final hashes and 97 hashes fixed before measurements\strut
\end{minipage}\tabularnewline
\begin{minipage}[t]{0.47\columnwidth}\raggedright
Independent numerical checks\strut
\end{minipage} & \begin{minipage}[t]{0.47\columnwidth}\raggedright
Reference outputs, guards, unchanged inputs, and state write-back
pass\strut
\end{minipage}\tabularnewline
\begin{minipage}[t]{0.47\columnwidth}\raggedright
Calibration and reserved tests\strut
\end{minipage} & \begin{minipage}[t]{0.47\columnwidth}\raggedright
Not measured\strut
\end{minipage}\tabularnewline
\bottomrule
\end{longtable}

The IME data call the actual M4 kernel with K32 packing. They exclude
production packing, tile scheduling, and backend synchronization. Four
threads retain the work per thread and increase total output columns.
The SSM data measure the complete convolution graph, including layout
conversion and state write-back. They keep total work fixed between
one-thread and four-thread settings. The convolution kernel length is 4
and the sequence count is 1. Both families reuse warm inputs from fixed
numerical fixtures. These data do not represent the shape frequencies or
cache states of full model inference.

The next table first takes the median paired time reduction for each
configuration. It then reports the median and range across the eight
configurations in that row. Positive values mean a faster candidate.
These ranges are observed training results, not prediction confidence
intervals.

\begin{longtable}[]{@{}lrrrl@{}}
\toprule
Operator scope & Threads & Configurations & Median reduction &
Configuration range\tabularnewline
\midrule
\endhead
IME M4 kernel slice & 1 & 8 & -12.90\% & -14.63\% to
-10.48\%\tabularnewline
IME M4 kernel slice & 4 & 8 & -12.73\% & -19.26\% to
-1.60\%\tabularnewline
SSM convolution graph & 1 & 8 & 37.68\% & 2.36\% to
92.00\%\tabularnewline
SSM convolution graph & 4 & 8 & -0.90\% & -41.26\% to
82.68\%\tabularnewline
\bottomrule
\end{longtable}

All 16 IME configurations have negative median reductions. Convolution
results depend on the shape and thread count. The batch retains all
candidates and all observations. It records 1,853 involuntary IME thread
switches. The generic instruction event's user filter still needs an
independent check. Foreground-thread cycles do not represent all
convolution graph threads.

The first batch chooses the common iteration count from the control arm
only. Its shortest formal sample is 0.864482 ms. Fixed overhead and
scheduler variation can affect short samples. The revised collector
probes both versions twice and raises the common count until all four
probes reach 10 ms. This revision has no verified second native batch.
Each operator family reserves 19 calibration shapes and nine independent
test shapes. Thread configurations of one shape remain in one group. The
model must freeze before calibration; its error bounds must freeze
before independent tests. Training observations cannot serve as
independent test evidence.

\hypertarget{connection-interruption-and-recovery-status}{%
\subsection{8.3 Connection interruption and recovery
status}\label{connection-interruption-and-recovery-status}}

K1 remote access fails during preparation for the second batch. The
available records cannot distinguish power loss, reboot, network
failure, or remote-service failure. They do not confirm a second-batch
start. A client disconnect does not establish that a board task stops.
The
\href{https://github.com/MoyaMryia/riscv-accl/blob/main/spacemit/reports/2026-10-12-k1-connection-interruption.md}{connection
report} retains that uncertainty.

The bounded remote client disables connection reuse and limits the total
client wait. It records unverified board task status as
\nolinkurl{unknown} and does not dispatch automatic retries. Six local
fake-client tests pass; actual recovery tests remain pending. Board work
waits for the requested physical check. After access returns, task,
tmux, lock, and batch records must be checked before any authorized
restart. This documentation update does not resume the stopped
combination or operator collection.

\hypertarget{reproduction-and-deliverables}{%
\section{9. Reproduction and
deliverables}\label{reproduction-and-deliverables}}

Use the
\href{https://github.com/MoyaMryia/riscv-accl/blob/main/spacemit/README.md}{deployment
guide} for model preparation and the historical patch stack. The
\href{https://github.com/MoyaMryia/riscv-accl/blob/main/spacemit/release/README.md}{current
infrastructure release} pins the later \nolinkurl{a990751} source and
packages its measured RVV/routing code with exact source hashes. The
older patch helper targets \nolinkurl{5ad05d8}; use separate checkouts.
Routing remains opt-in, and unsuccessful SSM/K32 candidates stay
disabled. The
\href{https://github.com/MoyaMryia/riscv-accl/blob/main/spacemit/reports/2026-10-07-release-verification.md}{fresh
release verification} passes a complete clean build and 104 exact
production graph cases per routing arm. All four code responses stop
naturally and pass 106 held-out checks each. All 35 archived artifacts
and source/build/runtime checks pass; both exit statuses are zero. The
one-pair natural-code observations show 508/521 tokens with identical
baseline/optimized/prior text. They support package verification,
without adding a new statistical speedup or general quality claim.

From the repository root, the latest checks are:

\begin{Shaded}
\begin{Highlighting}[]
\ExtensionTok{python3}\NormalTok{ spacemit/bench/test{-}k1{-}gemm{-}routing.py}
\ExtensionTok{python3}\NormalTok{ spacemit/bench/test{-}k1{-}fast{-}test.py}
\ExtensionTok{python3}\NormalTok{ spacemit/bench/test{-}k1{-}ime{-}test.py}
\ExtensionTok{python3}\NormalTok{ spacemit/bench/verify{-}k1{-}gemm{-}routing.py }\KeywordTok{\textbackslash{}}
  \ExtensionTok{spacemit/reports/raw/k1{-}gemm{-}routing{-}20261003{-}005740}
\end{Highlighting}
\end{Shaded}

To reproduce the board experiment, inspect its
\href{https://github.com/MoyaMryia/riscv-accl/blob/main/spacemit/experiments/2026-10-03-gemm-routing.md}{protocol}
and run \nolinkurl{python3 spacemit/bench/start-k1-gemm-routing.py} with
the matching board source/runtime/model prerequisites. The launcher runs
the benchmark and collector in tmux under a shared board lock, with
explicit time limits. The
\href{https://github.com/MoyaMryia/riscv-accl/blob/main/spacemit/bench/README.md}{benchmark
guide} covers the other campaigns.

The
\href{https://github.com/MoyaMryia/riscv-accl/blob/main/spacemit/reports/SUBMISSION-EVIDENCE.md}{evidence
appendix} maps every principal claim to source code, provenance, raw
records and verification commands. No model weights or private API
credentials are included in the report.

\hypertarget{release-decisions-and-optional-further-work}{%
\section{10. Release decisions and optional further
work}\label{release-decisions-and-optional-further-work}}

The bounded infrastructure release is verified and reproducible. The
current submission report, evidence appendix, source package and test
records form the completed deliverable. ARM/x86 comparison is excluded
from the release scope.

The strongest replicated long-input result is the activated RVV F16
attention kernel. The opt-in production GEMM routing also improves
measured prefill and bounded decode on both required models while
preserving tested outputs. Its scope is the current board/runtime and
the recorded workload; the release packages the measured mechanisms
within that scope. Routing remains opt-in.

Combined mode has independent 8k phase measurements and a
complete-answer pilot. The release uses direct decoding and keeps
routing opt-in. Baseline citation errors, native RS rollback mismatches
and the advanced context cases remain documented limitations; they
require separately declared follow-up protocols if pursued. Further
unused infrastructure leads are shared K packing across grouped query
heads, a Q4\_0 RVV lookup-table prototype, and compiler comparisons for
measured hotspots. Their platform fit and primary research references
are in the
\href{https://github.com/MoyaMryia/riscv-accl/blob/main/spacemit/reports/2026-10-03-infrastructure-research.md}{October
3 research review}. The
\href{https://github.com/MoyaMryia/riscv-accl/blob/main/spacemit/reports/2026-10-06-next-infrastructure-methods.md}{October
6 source audit} identifies channels-major SSM convolution with an RVV
channel loop and shared dense FFN activation packing. It also requires
separate decode attribution before changes to the worker schedule. The
\href{https://github.com/MoyaMryia/riscv-accl/blob/main/spacemit/reports/2026-10-06-ssm-conv-repaired-results.md}{repaired
convolution screen} now passes 432 exact graph cases in each of four
arms. RVV graph time falls 73.72\%/78.93\% at 32 tokens for the 2B/4B
shapes but rises 69.00\%/85.10\% at one token. This rejects
unconditional replacement; full-model state and timing were skipped in
that unconditional screen. The subsequent
\href{https://github.com/MoyaMryia/riscv-accl/blob/main/spacemit/reports/2026-10-06-ssm-complete-answer-results.md}{hybrid
complete-answer pilot} selects RVV only for batches of at least 32
tokens. Three arms pass 432 graph cases each; both models pass 16 exact
mixed-chunk/reset/RS-reference-fallback comparisons. All 24 answers stop
naturally and all 12 pairs are text-identical. Total answer latency
falls 2.92\%/2.13\% for 2B/4B, while absolute usefulness remains 2/6 and
5/6. Shared model errors and gains below the \textgreater3\% gate
prevent adoption. RS reference fallback does not resolve rollback
qualification. No production defaults change. The
\href{https://github.com/MoyaMryia/riscv-accl/blob/main/spacemit/reports/2026-10-06-local-clean-reference-results.md}{same-seed
clean local reference} uses the exact GGUFs and unpatched matching
source on generic x86 CPU, without custom K1 paths, flash attention or
MTP. All model/source/request hashes pass. It reproduces the 2B wrong
sum and missing citations, and both models still fail interval code.
Useful answers are 1/6 and 5/6 locally, versus 2/6 and 5/6 on both board
arms; one local 4B answer is length-capped. This supports limits of the
tested weights/settings without proving causality for every earlier
optimization across different CPUs. Local timing is not a matched x86/K1
performance comparison.

The
\href{https://github.com/MoyaMryia/riscv-accl/blob/main/spacemit/reports/2026-10-07-decode-packing-results.md}{October
7 sustained-decode diagnostic} passes 104 production graph cases in each
of three arms and eight matched cold requests with exact
prompt/token/text identity. Four decode profiles and all 86 artifact
hashes are verified. Duplicate FFN branch packing costs only
0.024--0.035\% of summed GEMM worker elapsed, so a shared-packing
implementation is not pursued. IME assembly accounts for 48.61--66.57\%
of sampled self CPU; approximately 20\% is calling-thread synchronous
waiting, which is not proven removable wall time. No new acceleration or
useful-answer result is claimed. The
\href{https://github.com/MoyaMryia/riscv-accl/blob/main/spacemit/reports/2026-10-07-ime-m1-k32-results.md}{fixed
K32 M1 inner-loop specialization} then passed 480 raw cases, 104
production cases per arm and all actual FFN/full output-head
comparisons. Its 144 operator measurements produced no gain above both
3\% and control spread; the candidate remains disabled. Full-model and
useful-answer stages were skipped, so it adds no measured inference
speedup. Separate load scheduling/dataflow, output-head partitioning and
dependency- preserving worker-barrier changes remain unmeasured leads.

Remaining advanced coverage limits include unmeasured 4B/32k and 64k
workloads, unresolved MTP long-code identity, and unavailable measured
GPU offload/full paged scheduling. These are stated explicitly so the
submission represents the completed work.

Validation priorities and primary-source practices for these remaining
checks are documented in the
\href{https://github.com/MoyaMryia/riscv-accl/blob/main/spacemit/reports/2026-10-03-validation-practice.md}{October
3 validation review}.

\Needspace{250pt}

\hypertarget{current-handover-and-remaining-work}{%
\section{11. Current handover and remaining
work}\label{current-handover-and-remaining-work}}

\begin{longtable}[]{@{}lll@{}}
\toprule
\begin{minipage}[b]{0.30\columnwidth}\raggedright
Area\strut
\end{minipage} & \begin{minipage}[b]{0.30\columnwidth}\raggedright
State at the evidence cutoff\strut
\end{minipage} & \begin{minipage}[b]{0.30\columnwidth}\raggedright
Next evidence needed\strut
\end{minipage}\tabularnewline
\midrule
\endhead
\begin{minipage}[t]{0.30\columnwidth}\raggedright
Pinned RVV/routing release\strut
\end{minipage} & \begin{minipage}[t]{0.30\columnwidth}\raggedright
Verified; routing stays opt-in\strut
\end{minipage} & \begin{minipage}[t]{0.30\columnwidth}\raggedright
A new protocol for any broader quality claim\strut
\end{minipage}\tabularnewline
\begin{minipage}[t]{0.30\columnwidth}\raggedright
Independent B1-r2/C1/C2-r2/C3 candidates\strut
\end{minipage} & \begin{minipage}[t]{0.30\columnwidth}\raggedright
Measured within their stated limits; not production defaults\strut
\end{minipage} & \begin{minipage}[t]{0.30\columnwidth}\raggedright
New release tests before integration\strut
\end{minipage}\tabularnewline
\begin{minipage}[t]{0.30\columnwidth}\raggedright
Combined-r1\strut
\end{minipage} & \begin{minipage}[t]{0.30\columnwidth}\raggedright
Failed one declared decode gate; stopped\strut
\end{minipage} & \begin{minipage}[t]{0.30\columnwidth}\raggedright
A separately declared revised-candidate protocol\strut
\end{minipage}\tabularnewline
\begin{minipage}[t]{0.30\columnwidth}\raggedright
Dispatch revision\strut
\end{minipage} & \begin{minipage}[t]{0.30\columnwidth}\raggedright
Local path checks and syntax check pass\strut
\end{minipage} & \begin{minipage}[t]{0.30\columnwidth}\raggedright
Native numerical, assembly, operator, and model checks\strut
\end{minipage}\tabularnewline
\begin{minipage}[t]{0.30\columnwidth}\raggedright
V7 evaluator\strut
\end{minipage} & \begin{minipage}[t]{0.30\columnwidth}\raggedright
Portable evidence replay passes; point directions unreliable\strut
\end{minipage} & \begin{minipage}[t]{0.30\columnwidth}\raggedright
Target-algorithm costs and independent error data\strut
\end{minipage}\tabularnewline
\begin{minipage}[t]{0.30\columnwidth}\raggedright
Operator data\strut
\end{minipage} & \begin{minipage}[t]{0.30\columnwidth}\raggedright
First training batch verified on two hosts\strut
\end{minipage} & \begin{minipage}[t]{0.30\columnwidth}\raggedright
Native checks for longer samples; separate calibration and tests\strut
\end{minipage}\tabularnewline
\begin{minipage}[t]{0.30\columnwidth}\raggedright
Board access\strut
\end{minipage} & \begin{minipage}[t]{0.30\columnwidth}\raggedright
Remote failure confirmed; physical cause unknown\strut
\end{minipage} & \begin{minipage}[t]{0.30\columnwidth}\raggedright
Physical check and bounded task/lock inspection\strut
\end{minipage}\tabularnewline
\begin{minipage}[t]{0.30\columnwidth}\raggedright
Initial-to-final speedup\strut
\end{minipage} & \begin{minipage}[t]{0.30\columnwidth}\raggedright
No matched complete comparison\strut
\end{minipage} & \begin{minipage}[t]{0.30\columnwidth}\raggedright
Fixed-source, fixed-model direct regression\strut
\end{minipage}\tabularnewline
\bottomrule
\end{longtable}

The strongest replicated long-input result remains the RVV attention
comparison. Routing adds measured phase gains on its optimized control.
Later individual gains do not override the combined regression or shared
quality failures. The project retains negative results, unresolved
mechanisms, and incomplete measurements.

\clearpage

\hypertarget{evidence-appendix}{%
\section{Evidence appendix}\label{evidence-appendix}}

Evidence cutoff: October 12, 2026, Asia/Singapore. This appendix
supports the
\href{https://github.com/MoyaMryia/riscv-accl/blob/main/spacemit/reports/SUBMISSION-REPORT.md}{submission
report}. File dates identify campaigns; historical work queues do not
describe current board activity.

\hypertarget{claim-to-artifact-map}{%
\subsection{1. Claim-to-artifact map}\label{claim-to-artifact-map}}

\begin{longtable}[]{@{}lll@{}}
\toprule
\begin{minipage}[b]{0.30\columnwidth}\raggedright
Submission claim\strut
\end{minipage} & \begin{minipage}[b]{0.30\columnwidth}\raggedright
Authoritative records\strut
\end{minipage} & \begin{minipage}[b]{0.30\columnwidth}\raggedright
Verification / implementation\strut
\end{minipage}\tabularnewline
\midrule
\endhead
\begin{minipage}[t]{0.30\columnwidth}\raggedright
RVV 2k/8k TTFT improvement on both models\strut
\end{minipage} & \begin{minipage}[t]{0.30\columnwidth}\raggedright
\href{https://github.com/MoyaMryia/riscv-accl/blob/main/spacemit/reports/raw/2026-09-25-lifecycle/lifecycle-rvv32.jsonl}{2k
requests},
\href{https://github.com/MoyaMryia/riscv-accl/blob/main/spacemit/reports/raw/2026-09-25-lifecycle/lifecycle-rvv32-long.jsonl}{8k
requests}\strut
\end{minipage} & \begin{minipage}[t]{0.30\columnwidth}\raggedright
\href{https://github.com/MoyaMryia/riscv-accl/blob/main/spacemit/bench/verify-results-matrix.py}{Matrix
verifier},
\href{https://github.com/MoyaMryia/riscv-accl/blob/main/spacemit/patches/0009-wide-rvv-vlen256.patch}{RVV
patch 0009}\strut
\end{minipage}\tabularnewline
\begin{minipage}[t]{0.30\columnwidth}\raggedright
Actual 16k pairs and 2B 32k feasibility\strut
\end{minipage} & \begin{minipage}[t]{0.30\columnwidth}\raggedright
\href{https://github.com/MoyaMryia/riscv-accl/blob/main/spacemit/reports/2026-09-27-resumed-measures.md}{Resumed
measurement matrix},
\href{https://github.com/MoyaMryia/riscv-accl/blob/main/spacemit/reports/raw/2026-09-25-lifecycle}{lifecycle
raw directory}\strut
\end{minipage} & \begin{minipage}[t]{0.30\columnwidth}\raggedright
Same matrix verifier; individual server logs and requests are linked in
the matrix report\strut
\end{minipage}\tabularnewline
\begin{minipage}[t]{0.30\columnwidth}\raggedright
New routing model results\strut
\end{minipage} & \begin{minipage}[t]{0.30\columnwidth}\raggedright
\href{https://github.com/MoyaMryia/riscv-accl/blob/main/spacemit/reports/raw/k1-gemm-routing-20261003-005740/summary.json}{Routing
summary},
\href{https://github.com/MoyaMryia/riscv-accl/blob/main/spacemit/reports/raw/k1-gemm-routing-20261003-005740}{individual
requests/logs}\strut
\end{minipage} & \begin{minipage}[t]{0.30\columnwidth}\raggedright
\href{https://github.com/MoyaMryia/riscv-accl/blob/main/spacemit/bench/verify-k1-gemm-routing.py}{Independent
verifier},
\href{https://github.com/MoyaMryia/riscv-accl/blob/main/spacemit/reports/raw/k1-gemm-routing-20261003-005740/verification.json}{verification
output}\strut
\end{minipage}\tabularnewline
\begin{minipage}[t]{0.30\columnwidth}\raggedright
Combined routing at 8k, complete-answer pilot and RS replay
failures\strut
\end{minipage} & \begin{minipage}[t]{0.30\columnwidth}\raggedright
\href{https://github.com/MoyaMryia/riscv-accl/blob/main/spacemit/reports/2026-10-03-staged-validation-results.md}{Staged
validation results},
\href{https://github.com/MoyaMryia/riscv-accl/blob/main/spacemit/reports/raw/k1-validation-20261003-124720/summary.json}{summary},
\href{https://github.com/MoyaMryia/riscv-accl/blob/main/spacemit/reports/raw/k1-validation-20261003-124720/collection-receipt.json}{83-artifact
receipt}\strut
\end{minipage} & \begin{minipage}[t]{0.30\columnwidth}\raggedright
\href{https://github.com/MoyaMryia/riscv-accl/blob/main/spacemit/bench/verify-k1-validation.py}{Validation
verifier},
\href{https://github.com/MoyaMryia/riscv-accl/blob/main/spacemit/reports/raw/k1-validation-20261003-124720/verification.json}{verification
output}; 4B/32k explicitly skipped\strut
\end{minipage}\tabularnewline
\begin{minipage}[t]{0.30\columnwidth}\raggedright
Native and operator routing gates\strut
\end{minipage} & \begin{minipage}[t]{0.30\columnwidth}\raggedright
\href{https://github.com/MoyaMryia/riscv-accl/blob/main/spacemit/reports/raw/k1-gemm-routing-20261003-005740/operator.jsonl}{Operator
records}, numerical logs in the same directory\strut
\end{minipage} & \begin{minipage}[t]{0.30\columnwidth}\raggedright
\href{https://github.com/MoyaMryia/riscv-accl/blob/main/spacemit/bench/test-k1-gemm-routing.cpp}{Production
graph harness},
\href{https://github.com/MoyaMryia/riscv-accl/blob/main/spacemit/bench/k1-gemm-routing.py}{controller}\strut
\end{minipage}\tabularnewline
\begin{minipage}[t]{0.30\columnwidth}\raggedright
236 collected artifacts and baseline preservation\strut
\end{minipage} & \begin{minipage}[t]{0.30\columnwidth}\raggedright
\href{https://github.com/MoyaMryia/riscv-accl/blob/main/spacemit/reports/raw/k1-gemm-routing-20261003-005740/collection-receipt.json}{Collection
receipt},
\href{https://github.com/MoyaMryia/riscv-accl/blob/main/spacemit/reports/raw/k1-gemm-routing-20261003-005740/baseline-preservation-check.json}{46-file
preservation check}\strut
\end{minipage} & \begin{minipage}[t]{0.30\columnwidth}\raggedright
Independent artifact hashing; preservation compares end-of-run board
hashes to starting manifests\strut
\end{minipage}\tabularnewline
\begin{minipage}[t]{0.30\columnwidth}\raggedright
Source/build/model provenance\strut
\end{minipage} & \begin{minipage}[t]{0.30\columnwidth}\raggedright
\href{https://github.com/MoyaMryia/riscv-accl/blob/main/spacemit/reports/raw/k1-gemm-routing-20261003-005740/expected-provenance.json}{Expected
baseline},
\href{https://github.com/MoyaMryia/riscv-accl/blob/main/spacemit/reports/raw/k1-gemm-routing-20261003-005740/provenance.json}{build
provenance},
\href{https://github.com/MoyaMryia/riscv-accl/blob/main/spacemit/reports/raw/k1-gemm-routing-20261003-005740/candidate-gemm-routing.patch}{tested
patch}\strut
\end{minipage} & \begin{minipage}[t]{0.30\columnwidth}\raggedright
\href{https://github.com/MoyaMryia/riscv-accl/blob/main/spacemit/bench/make-k1-gemm-routing.py}{Generator},
\href{https://github.com/MoyaMryia/riscv-accl/blob/main/spacemit/bench/start-k1-gemm-routing.py}{launcher}\strut
\end{minipage}\tabularnewline
\begin{minipage}[t]{0.30\columnwidth}\raggedright
Q8\_0 IME1 and M4 scale results\strut
\end{minipage} & \begin{minipage}[t]{0.30\columnwidth}\raggedright
\href{https://github.com/MoyaMryia/riscv-accl/blob/main/spacemit/reports/2026-09-24-integrated-k1.md}{September
24 integrated report},
\href{https://github.com/MoyaMryia/riscv-accl/blob/main/spacemit/reports/raw/2026-09-24-integrated}{raw
arms}\strut
\end{minipage} & \begin{minipage}[t]{0.30\columnwidth}\raggedright
\href{https://github.com/MoyaMryia/riscv-accl/blob/main/spacemit/patches/0007-q8-ime1.patch}{Patch
0007},
\href{https://github.com/MoyaMryia/riscv-accl/blob/main/spacemit/patches/0008-ime-m4-scale.patch}{patch
0008}\strut
\end{minipage}\tabularnewline
\begin{minipage}[t]{0.30\columnwidth}\raggedright
Corrected small-effect statistics\strut
\end{minipage} & \begin{minipage}[t]{0.30\columnwidth}\raggedright
\href{https://github.com/MoyaMryia/riscv-accl/blob/main/spacemit/reports/2026-09-28-local-verification.md}{Local
verification}\strut
\end{minipage} & \begin{minipage}[t]{0.30\columnwidth}\raggedright
\href{https://github.com/MoyaMryia/riscv-accl/blob/main/spacemit/bench/paired-stats.py}{Statistics
script}; fractional Welch degrees of freedom\strut
\end{minipage}\tabularnewline
\begin{minipage}[t]{0.30\columnwidth}\raggedright
MTP long-code mismatch and bounded decode traces\strut
\end{minipage} & \begin{minipage}[t]{0.30\columnwidth}\raggedright
\href{https://github.com/MoyaMryia/riscv-accl/blob/main/spacemit/reports/2026-09-27-resumed-measures.md}{Resumed
diagnostics},
\href{https://github.com/MoyaMryia/riscv-accl/blob/main/spacemit/reports/raw/2026-09-25-lifecycle/lifecycle-code-long.jsonl}{long-code
traces}\strut
\end{minipage} & \begin{minipage}[t]{0.30\columnwidth}\raggedright
\href{https://github.com/MoyaMryia/riscv-accl/blob/main/spacemit/bench/audit-lifecycle.py}{Lifecycle
auditor}; smaller drafts/RS/target-logit/FA-off evidence in the resumed
report\strut
\end{minipage}\tabularnewline
\begin{minipage}[t]{0.30\columnwidth}\raggedright
Complete-answer checkpoint MTP usefulness\strut
\end{minipage} & \begin{minipage}[t]{0.30\columnwidth}\raggedright
\href{https://github.com/MoyaMryia/riscv-accl/blob/main/spacemit/reports/2026-10-03-mtp-usefulness-results.md}{Results},
\href{https://github.com/MoyaMryia/riscv-accl/blob/main/spacemit/reports/raw/mtp-quality-20261003-173915/quality-summary.json}{quality
summary},
\href{https://github.com/MoyaMryia/riscv-accl/blob/main/spacemit/reports/raw/mtp-quality-20261003-173915/collection-receipt.json}{21-artifact
receipt}\strut
\end{minipage} & \begin{minipage}[t]{0.30\columnwidth}\raggedright
\href{https://github.com/MoyaMryia/riscv-accl/blob/main/spacemit/experiments/2026-10-03-mtp-usefulness.md}{Protocol},
\href{https://github.com/MoyaMryia/riscv-accl/blob/main/spacemit/bench/mtp-quality-checks.py}{functional
checks},
\href{https://github.com/MoyaMryia/riscv-accl/blob/main/spacemit/reports/raw/mtp-quality-20261003-173915/scoring-recovery.json}{scoring
recovery}; relative gates pass, absolute usefulness fails\strut
\end{minipage}\tabularnewline
\begin{minipage}[t]{0.30\columnwidth}\raggedright
Adaptive MTP bounded infrastructure screen\strut
\end{minipage} & \begin{minipage}[t]{0.30\columnwidth}\raggedright
\href{https://github.com/MoyaMryia/riscv-accl/blob/main/spacemit/reports/2026-10-05-adaptive-mtp-results.md}{Results},
\href{https://github.com/MoyaMryia/riscv-accl/blob/main/spacemit/reports/raw/adaptive-mtp-infra-20261005-153716/adaptive-quality-summary.json}{final
summary},
\href{https://github.com/MoyaMryia/riscv-accl/blob/main/spacemit/reports/raw/adaptive-mtp-infra-20261005-153716/collection-receipt.json}{47-artifact
receipt}\strut
\end{minipage} & \begin{minipage}[t]{0.30\columnwidth}\raggedright
\href{https://github.com/MoyaMryia/riscv-accl/blob/main/spacemit/experiments/2026-10-05-adaptive-mtp-infrastructure.md}{Protocol},
\href{https://github.com/MoyaMryia/riscv-accl/blob/main/spacemit/reports/raw/adaptive-mtp-infra-20261005-153716/adaptive-scoring-recovery.json}{grading
recovery}; throughput positive, overall candidate unqualified\strut
\end{minipage}\tabularnewline
\begin{minipage}[t]{0.30\columnwidth}\raggedright
Complete-answer cache matrix\strut
\end{minipage} & \begin{minipage}[t]{0.30\columnwidth}\raggedright
\href{https://github.com/MoyaMryia/riscv-accl/blob/main/spacemit/reports/2026-09-30-chat-quality-benchmark.md}{Quality
report},
\href{https://github.com/MoyaMryia/riscv-accl/blob/main/spacemit/reports/raw/document-quality-20260930-full-v2/summary.md}{completed
24-pair summary}\strut
\end{minipage} & \begin{minipage}[t]{0.30\columnwidth}\raggedright
\href{https://github.com/MoyaMryia/riscv-accl/blob/main/spacemit/bench/run-all-document-quality.py}{Quality
runner},
\href{https://github.com/MoyaMryia/riscv-accl/blob/main/spacemit/bench/bench-shared-document-cache.py}{document
benchmark}\strut
\end{minipage}\tabularnewline
\begin{minipage}[t]{0.30\columnwidth}\raggedright
Negative compact-layout screen\strut
\end{minipage} & \begin{minipage}[t]{0.30\columnwidth}\raggedright
\href{https://github.com/MoyaMryia/riscv-accl/blob/main/spacemit/reports/2026-10-01-fast-method.md}{Fast
method results},
\href{https://github.com/MoyaMryia/riscv-accl/blob/main/spacemit/reports/raw/k1-fast-20261001-224957/summary.md}{completed
run}\strut
\end{minipage} & \begin{minipage}[t]{0.30\columnwidth}\raggedright
\href{https://github.com/MoyaMryia/riscv-accl/blob/main/spacemit/experiments/2026-09-30-fast-test-design.md}{Staged
test design}\strut
\end{minipage}\tabularnewline
\begin{minipage}[t]{0.30\columnwidth}\raggedright
Negative IME / recurrent / attention screens\strut
\end{minipage} & \begin{minipage}[t]{0.30\columnwidth}\raggedright
\href{https://github.com/MoyaMryia/riscv-accl/blob/main/spacemit/reports/2026-10-02-ime-scheduling.md}{IME
scheduling},
\href{https://github.com/MoyaMryia/riscv-accl/blob/main/spacemit/reports/2026-10-02-recurrent-prefill.md}{recurrent
fusion},
\href{https://github.com/MoyaMryia/riscv-accl/blob/main/spacemit/experiments/2026-10-02-attention-infrastructure.md}{attention
infrastructure}\strut
\end{minipage} & \begin{minipage}[t]{0.30\columnwidth}\raggedright
Their linked native tests, operator records and gated model runs\strut
\end{minipage}\tabularnewline
\begin{minipage}[t]{0.30\columnwidth}\raggedright
GPU and page-gather boundaries\strut
\end{minipage} & \begin{minipage}[t]{0.30\columnwidth}\raggedright
\href{https://github.com/MoyaMryia/riscv-accl/blob/main/spacemit/reports/2026-09-27-completed-gates.md}{Completed
gates},
\href{https://github.com/MoyaMryia/riscv-accl/blob/main/spacemit/reports/2026-09-27-resumed-measures.md}{later
telemetry and backend findings}\strut
\end{minipage} & \begin{minipage}[t]{0.30\columnwidth}\raggedright
Zero transferred OpenCL layers; Vulkan dropped; page gathering retains
full backing allocation\strut
\end{minipage}\tabularnewline
\begin{minipage}[t]{0.30\columnwidth}\raggedright
Independent bandwidth and routing confirmation\strut
\end{minipage} & \begin{minipage}[t]{0.30\columnwidth}\raggedright
\href{https://github.com/MoyaMryia/riscv-accl/blob/main/spacemit/reports/2026-10-06-k1-roofline-results.md}{October
6 results},
\href{https://github.com/MoyaMryia/riscv-accl/blob/main/spacemit/reports/raw/k1-roofline-20261005-225502/summary.json}{summary},
\href{https://github.com/MoyaMryia/riscv-accl/blob/main/spacemit/reports/raw/k1-roofline-20261005-225502/collection-receipt.json}{150-artifact
receipt}\strut
\end{minipage} & \begin{minipage}[t]{0.30\columnwidth}\raggedright
\href{https://github.com/MoyaMryia/riscv-accl/blob/main/spacemit/experiments/2026-10-05-k1-roofline.md}{Protocol},
\href{https://github.com/MoyaMryia/riscv-accl/blob/main/spacemit/reports/raw/k1-roofline-20261005-225502/ARCHIVE.md}{compressed
evidence storage}; traffic references are not hardware-limit proof\strut
\end{minipage}\tabularnewline
\begin{minipage}[t]{0.30\columnwidth}\raggedright
Same-seed local clean-reference quality\strut
\end{minipage} & \begin{minipage}[t]{0.30\columnwidth}\raggedright
\href{https://github.com/MoyaMryia/riscv-accl/blob/main/spacemit/reports/2026-10-06-local-clean-reference-results.md}{Results},
\href{https://github.com/MoyaMryia/riscv-accl/blob/main/spacemit/reports/raw/local-reference-quality-20261006-194015/summary.json}{summary},
\href{https://github.com/MoyaMryia/riscv-accl/blob/main/spacemit/reports/raw/local-reference-quality-20261006-194015/completion-verification.json}{verification}\strut
\end{minipage} & \begin{minipage}[t]{0.30\columnwidth}\raggedright
\href{https://github.com/MoyaMryia/riscv-accl/blob/main/spacemit/experiments/2026-10-06-local-clean-reference.md}{Protocol},
exact GGUF/source/request hashes, 12 audits and 11 two-order judgment
pairs; errors persist, cross-backend causal/speed claims excluded\strut
\end{minipage}\tabularnewline
\begin{minipage}[t]{0.30\columnwidth}\raggedright
Hybrid SSM complete useful answers\strut
\end{minipage} & \begin{minipage}[t]{0.30\columnwidth}\raggedright
\href{https://github.com/MoyaMryia/riscv-accl/blob/main/spacemit/reports/2026-10-06-ssm-complete-answer-results.md}{Measured
results},
\href{https://github.com/MoyaMryia/riscv-accl/blob/main/spacemit/reports/raw/k1-ssm-quality-20261006-131234/ssm-quality-summary.json}{final
quality summary},
\href{https://github.com/MoyaMryia/riscv-accl/blob/main/spacemit/reports/raw/k1-ssm-quality-20261006-131234/collection-receipt.json}{114-artifact
receipt}\strut
\end{minipage} & \begin{minipage}[t]{0.30\columnwidth}\raggedright
\href{https://github.com/MoyaMryia/riscv-accl/blob/main/spacemit/reports/raw/k1-ssm-quality-20261006-131234/completion-verification.json}{Completion
verification},
\href{https://github.com/MoyaMryia/riscv-accl/blob/main/spacemit/reports/2026-10-06-ssm-useful-answer-example.md}{passing
answer}; exact graph/state pass, 24 natural answers, adoption gates
fail; local scoring files additional to receipt\strut
\end{minipage}\tabularnewline
\begin{minipage}[t]{0.30\columnwidth}\raggedright
Sustained decode and negative shared-packing priority\strut
\end{minipage} & \begin{minipage}[t]{0.30\columnwidth}\raggedright
\href{https://github.com/MoyaMryia/riscv-accl/blob/main/spacemit/reports/2026-10-07-decode-packing-results.md}{October
7 results},
\href{https://github.com/MoyaMryia/riscv-accl/blob/main/spacemit/reports/raw/k1-decode-audit-20261007-074314/summary.json}{summary},
\href{https://github.com/MoyaMryia/riscv-accl/blob/main/spacemit/reports/raw/k1-decode-audit-20261007-074314/collection-receipt.json}{86-artifact
receipt}\strut
\end{minipage} & \begin{minipage}[t]{0.30\columnwidth}\raggedright
\href{https://github.com/MoyaMryia/riscv-accl/blob/main/spacemit/reports/raw/k1-decode-analysis-20261007/verify-completion.py}{Independent
verifier},
\href{https://github.com/MoyaMryia/riscv-accl/blob/main/spacemit/reports/raw/k1-decode-analysis-20261007/verification.json}{verification};
eight exact matched requests, four decode profiles, native identity;
packing declined, CPU shares are not wall-time savings\strut
\end{minipage}\tabularnewline
\begin{minipage}[t]{0.30\columnwidth}\raggedright
Unused research opportunities\strut
\end{minipage} & \begin{minipage}[t]{0.30\columnwidth}\raggedright
\href{https://github.com/MoyaMryia/riscv-accl/blob/main/spacemit/reports/2026-10-03-infrastructure-research.md}{October
3 GitHub/arXiv/Scholar review},
\href{https://github.com/MoyaMryia/riscv-accl/blob/main/spacemit/reports/2026-10-06-next-infrastructure-methods.md}{October
6 source audit}\strut
\end{minipage} & \begin{minipage}[t]{0.30\columnwidth}\raggedright
Dated primary-source/source reviews; convolution and decode packing
subsequently measured; M1 kernel scheduling and worker-barrier
implementation remain unqualified\strut
\end{minipage}\tabularnewline
\bottomrule
\end{longtable}

\hypertarget{baseline-identities}{%
\subsection{2. Baseline identities}\label{baseline-identities}}

\begin{longtable}[]{@{}ll@{}}
\toprule
\begin{minipage}[b]{0.47\columnwidth}\raggedright
Source state\strut
\end{minipage} & \begin{minipage}[b]{0.47\columnwidth}\raggedright
Role\strut
\end{minipage}\tabularnewline
\midrule
\endhead
\begin{minipage}[t]{0.47\columnwidth}\raggedright
Official SpaceMiT fork \nolinkurl{5ad05d8}\strut
\end{minipage} & \begin{minipage}[t]{0.47\columnwidth}\raggedright
Clean starting point for the packaged patch helper\strut
\end{minipage}\tabularnewline
\begin{minipage}[t]{0.47\columnwidth}\raggedright
\nolinkurl{6562c22} plus separately selected options\strut
\end{minipage} & \begin{minipage}[t]{0.47\columnwidth}\raggedright
Earlier recurrent/MTP integration state\strut
\end{minipage}\tabularnewline
\begin{minipage}[t]{0.47\columnwidth}\raggedright
\nolinkurl{a990751} plus recorded wide-RVV source changes\strut
\end{minipage} & \begin{minipage}[t]{0.47\columnwidth}\raggedright
Later board baseline; build-specific provenance is required\strut
\end{minipage}\tabularnewline
\begin{minipage}[t]{0.47\columnwidth}\raggedright
Original library versus generated mode 0\strut
\end{minipage} & \begin{minipage}[t]{0.47\columnwidth}\raggedright
Routing operator control for unintended effects of the isolated library
build\strut
\end{minipage}\tabularnewline
\begin{minipage}[t]{0.47\columnwidth}\raggedright
Generated mode 0 versus mode 1 or 2\strut
\end{minipage} & \begin{minipage}[t]{0.47\columnwidth}\raggedright
Routing model comparisons; the baseline already has wide
attention/layout 0\strut
\end{minipage}\tabularnewline
\bottomrule
\end{longtable}

Historical documentation values with different quantization, model
derivative, prompt length or runtime state are not matched baseline
arms. In particular, the old Q4\_1 documentation rates must not be used
to calculate a controlled Q4\_0 optimization percentage. Baselines from
different campaigns must not be treated as one continuous cumulative
curve.

\hypertarget{latest-model-identity}{%
\subsection{3. Latest model identity}\label{latest-model-identity}}

The routing run hashes the exact local files; model preparation is
described in the
\href{https://github.com/MoyaMryia/riscv-accl/blob/main/spacemit/models/tools/README.md}{model
tools guide}.

\begin{longtable}[]{@{}ll@{}}
\toprule
\begin{minipage}[b]{0.47\columnwidth}\raggedright
Model file\strut
\end{minipage} & \begin{minipage}[b]{0.47\columnwidth}\raggedright
SHA-256\strut
\end{minipage}\tabularnewline
\midrule
\endhead
\begin{minipage}[t]{0.47\columnwidth}\raggedright
\nolinkurl{Qwen3.5-2B-MTP-Q4_0-embQ4_0-dv64k.gguf}\strut
\end{minipage} & \begin{minipage}[t]{0.47\columnwidth}\raggedright
\texttt{56f0ddd9\allowbreak{}0dfa0e24\allowbreak{}56af6cb4\allowbreak{}718ece41\allowbreak{}9678d21a\allowbreak{}75d65e27\allowbreak{}353319b5\allowbreak{}e431a2fa}\strut
\end{minipage}\tabularnewline
\begin{minipage}[t]{0.47\columnwidth}\raggedright
\nolinkurl{Qwen3.5-4B-MTP-Q4_0-embQ4_0-dv64k.gguf}\strut
\end{minipage} & \begin{minipage}[t]{0.47\columnwidth}\raggedright
\texttt{0adf6cce\allowbreak{}5df53921\allowbreak{}606033c3\allowbreak{}9a13fae2\allowbreak{}054543b7\allowbreak{}8af0e7ff\allowbreak{}98860e86\allowbreak{}f078ac0f}\strut
\end{minipage}\tabularnewline
\bottomrule
\end{longtable}

Model weights remain outside this repository. The current kernel
comparisons do not replace the model, alter its Q4\_0 weights or trim
input tokens.

\hypertarget{re-run-local-verification}{%
\subsection{4. Re-run local
verification}\label{re-run-local-verification}}

From the repository root:

\begin{Shaded}
\begin{Highlighting}[]
\ExtensionTok{python3}\NormalTok{ spacemit/bench/verify{-}results{-}matrix.py }\KeywordTok{\textbackslash{}}
  \ExtensionTok{spacemit/reports/2026{-}09{-}27{-}resumed{-}measures.md} \KeywordTok{\textbackslash{}}
  \ExtensionTok{spacemit/reports/raw/2026{-}09{-}25{-}lifecycle}
\ExtensionTok{python3}\NormalTok{ spacemit/bench/verify{-}k1{-}gemm{-}routing.py }\KeywordTok{\textbackslash{}}
  \ExtensionTok{spacemit/reports/raw/k1{-}gemm{-}routing{-}20261003{-}005740}
\ExtensionTok{python3}\NormalTok{ spacemit/bench/test{-}k1{-}gemm{-}routing.py}
\ExtensionTok{python3}\NormalTok{ spacemit/bench/test{-}k1{-}fast{-}test.py}
\ExtensionTok{python3}\NormalTok{ spacemit/bench/test{-}k1{-}ime{-}test.py}
\end{Highlighting}
\end{Shaded}

The matrix verifier checks 21 rows / 105 numeric cells. The routing
verifier checks archived hashes, numerical logs, operator and model
gates, input/output counts, actual runtime configuration and per-phase
kernel markers. The three local regression scripts contain 2, 6 and 3
checks respectively. These local tests supplement the board runs; they
do not substitute for them.

The routing run's environment telemetry allowlist omitted
\nolinkurl{SPINE_K1_GEMM_ROUTE}. Actual mandatory kernel markers record
mode, eligibility and bypass for each phase and establish activation
independently. The verifier requires those markers and checks the
environment value if recorded. Future runs also capture the environment
field. Original raw records were preserved.

The 236-artifact receipt covers collected board evidence. The
independently generated local verification and preservation records are
additional files, not members of that original collection count. Native
binary output dumps remain on the board; numerical logs and summaries
record their identical SHA-256 values. End-of-run board hashing
independently confirmed those dumps.

\hypertarget{fixed-k32-m1-screen-negative-evidence}{%
\subsection{Fixed K32 M1 screen: negative
evidence}\label{fixed-k32-m1-screen-negative-evidence}}

The
\href{https://github.com/MoyaMryia/riscv-accl/blob/main/spacemit/reports/2026-10-07-ime-m1-k32-results.md}{October
7 M1 specialization results} have the following evidence: the
\href{https://github.com/MoyaMryia/riscv-accl/blob/main/spacemit/reports/raw/k1-ime-m1-k32-20261007-121231/summary.json}{frozen
summary},
\href{https://github.com/MoyaMryia/riscv-accl/blob/main/spacemit/reports/raw/k1-ime-m1-k32-20261007-121231/operator.jsonl}{144
operator samples},
\href{https://github.com/MoyaMryia/riscv-accl/blob/main/spacemit/reports/raw/k1-ime-m1-k32-20261007-121231/collection-receipt.json}{211-artifact
receipt},
\href{https://github.com/MoyaMryia/riscv-accl/blob/main/spacemit/reports/raw/k1-ime-m1-analysis-20261007/verify-completion.py}{independent
verifier},
\href{https://github.com/MoyaMryia/riscv-accl/blob/main/spacemit/reports/raw/k1-ime-m1-analysis-20261007/verification.json}{verification
output} and
\href{https://github.com/MoyaMryia/riscv-accl/blob/main/spacemit/reports/raw/k1-ime-m1-analysis-20261007/remote-preservation.json}{40-file
remote preservation check}. Both exit statuses are zero. All 480 raw
cases, 104 production cases per arm and six actual
FFN/full-248320-column shapes per arm pass exact outputs and guards.
Recreated route-3 baseline hash is exact; generated assembly removes the
inner counter/branches while retaining dot/load/FMA order and fallbacks.
No operator clears both 3\% and control spread. Full-model state, cold
ABBA and useful-answer stages were skipped. Candidate disabled; no new
inference or quality gain enters the submission. Independent analysis
files are outside the original 211-artifact receipt.

\hypertarget{submission-document-roles}{%
\subsection{5. Submission document
roles}\label{submission-document-roles}}

The
\href{https://github.com/MoyaMryia/riscv-accl/blob/main/spacemit/reports/2026-10-06-ssm-conv-repaired-results.md}{October
6 repaired SSM screen} is supported by
\href{https://github.com/MoyaMryia/riscv-accl/blob/main/spacemit/reports/raw/k1-ssm-conv-20261006-115459/summary.json}{its
summary},
\href{https://github.com/MoyaMryia/riscv-accl/blob/main/spacemit/reports/raw/k1-ssm-conv-20261006-115459/operator.jsonl}{operator
samples} and
\href{https://github.com/MoyaMryia/riscv-accl/blob/main/spacemit/reports/raw/k1-ssm-conv-20261006-115459/collection-receipt.json}{97-artifact
receipt}. All four arms pass 432 exact output/history cases. Warm graph
gains at 32 tokens accompany single-token regressions, so model-state
and inference timing stages were skipped. These measurements support no
full-model gain.

\begin{longtable}[]{@{}ll@{}}
\toprule
\begin{minipage}[b]{0.47\columnwidth}\raggedright
Material\strut
\end{minipage} & \begin{minipage}[b]{0.47\columnwidth}\raggedright
Use\strut
\end{minipage}\tabularnewline
\midrule
\endhead
\begin{minipage}[t]{0.47\columnwidth}\raggedright
\href{https://github.com/MoyaMryia/riscv-accl/blob/main/spacemit/reports/SUBMISSION-REPORT.md}{SUBMISSION-REPORT.md}\strut
\end{minipage} & \begin{minipage}[t]{0.47\columnwidth}\raggedright
Current report; edit this file before a new export\strut
\end{minipage}\tabularnewline
\begin{minipage}[t]{0.47\columnwidth}\raggedright
\href{https://github.com/MoyaMryia/riscv-accl/blob/main/spacemit/reports/SUBMISSION-REPORT.tex}{SUBMISSION-REPORT.tex}\strut
\end{minipage} & \begin{minipage}[t]{0.47\columnwidth}\raggedright
Standalone TeX export of the report and this appendix\strut
\end{minipage}\tabularnewline
\begin{minipage}[t]{0.47\columnwidth}\raggedright
\href{https://github.com/MoyaMryia/riscv-accl/blob/main/spacemit/reports/SUBMISSION-REPORT.pdf}{SUBMISSION-REPORT.pdf}\strut
\end{minipage} & \begin{minipage}[t]{0.47\columnwidth}\raggedright
PDF export for review\strut
\end{minipage}\tabularnewline
\begin{minipage}[t]{0.47\columnwidth}\raggedright
\href{https://github.com/MoyaMryia/riscv-accl/blob/main/spacemit/reports/REPORT-BUILD.md}{REPORT-BUILD.md}\strut
\end{minipage} & \begin{minipage}[t]{0.47\columnwidth}\raggedright
Export commands and document update rules\strut
\end{minipage}\tabularnewline
\begin{minipage}[t]{0.47\columnwidth}\raggedright
This appendix\strut
\end{minipage} & \begin{minipage}[t]{0.47\columnwidth}\raggedright
Evidence and reproduction references accompanying the report\strut
\end{minipage}\tabularnewline
\begin{minipage}[t]{0.47\columnwidth}\raggedright
\href{https://github.com/MoyaMryia/riscv-accl/blob/main/spacemit/DOCS.md}{Current
documentation guide}\strut
\end{minipage} & \begin{minipage}[t]{0.47\columnwidth}\raggedright
Operational status and navigation\strut
\end{minipage}\tabularnewline
\begin{minipage}[t]{0.47\columnwidth}\raggedright
\href{https://github.com/MoyaMryia/riscv-accl/blob/main/spacemit/README.md}{Integration
guide} /
\href{https://github.com/MoyaMryia/riscv-accl/blob/main/spacemit/bench/README.md}{benchmark
guide}\strut
\end{minipage} & \begin{minipage}[t]{0.47\columnwidth}\raggedright
Deployment and executable procedures\strut
\end{minipage}\tabularnewline
\begin{minipage}[t]{0.47\columnwidth}\raggedright
Dated reports and experiments\strut
\end{minipage} & \begin{minipage}[t]{0.47\columnwidth}\raggedright
Campaign-specific details and historical decisions\strut
\end{minipage}\tabularnewline
\begin{minipage}[t]{0.47\columnwidth}\raggedright
\nolinkurl{raw/}\strut
\end{minipage} & \begin{minipage}[t]{0.47\columnwidth}\raggedright
Reproducibility records; preserve original measurements\strut
\end{minipage}\tabularnewline
\begin{minipage}[t]{0.47\columnwidth}\raggedright
\href{https://github.com/MoyaMryia/riscv-accl/blob/main/spacemit/reports/k1-llm-inference-optimizations.tex}{September
24 LaTeX} and
\href{https://github.com/MoyaMryia/riscv-accl/blob/main/spacemit/reports/k1-llm-inference-optimizations.pdf}{PDF}\strut
\end{minipage} & \begin{minipage}[t]{0.47\columnwidth}\raggedright
Historical export; predates later findings and is not the current
submission\strut
\end{minipage}\tabularnewline
\begin{minipage}[t]{0.47\columnwidth}\raggedright
\href{https://github.com/MoyaMryia/riscv-accl/blob/main/FINAL-REPORT.md}{Root
FINAL-REPORT.md}\strut
\end{minipage} & \begin{minipage}[t]{0.47\columnwidth}\raggedright
Historical short-prompt campaign, with later annotations; not the
current submission\strut
\end{minipage}\tabularnewline
\bottomrule
\end{longtable}

No required final-report template, author list or external submission
destination has been supplied. The report is organized for review; it
has not been submitted externally.

\hypertarget{current-release-package}{%
\subsection{Current release package}\label{current-release-package}}

The
\href{https://github.com/MoyaMryia/riscv-accl/blob/main/spacemit/release/manifest.json}{release
manifest} pins the exact source revision, patch and measured files,
model hashes and matched route0/3 profiles.
\href{https://github.com/MoyaMryia/riscv-accl/blob/main/spacemit/reports/raw/release-20261007/clean-apply-verification.json}{Offline
clean-source verification} checks the 3,266-file snapshot and exact
patch round trip. The
\href{https://github.com/MoyaMryia/riscv-accl/blob/main/spacemit/reports/2026-10-07-release-verification.md}{fresh
full board build and complete-answer verification} passes in
\nolinkurl{k1-release-check-20261007-125712}. It checks 104 exact
production graph cases per routing arm. All four Unicode code answers
stop naturally and pass 106 held-out checks each. All 35 archived
artifacts, 18 staged code files and three collector files pass checksum
verification; post-run source, binary, compiler and vendor-runtime
checks pass. Board/collector exit statuses are zero. The
\href{https://github.com/MoyaMryia/riscv-accl/blob/main/spacemit/reports/raw/release-20261007/verification.json}{independent
result} records exact paired/prior text identity and the descriptive
natural-answer timings. ARM/x86 comparison is excluded by user
instruction. Release smoke timings do not add a new statistical speedup
claim.

\hypertarget{evidence-added-after-october-7}{%
\subsection{6. Evidence added after October
7}\label{evidence-added-after-october-7}}

\begin{longtable}[]{@{}lll@{}}
\toprule
\begin{minipage}[b]{0.30\columnwidth}\raggedright
Claim\strut
\end{minipage} & \begin{minipage}[b]{0.30\columnwidth}\raggedright
Records\strut
\end{minipage} & \begin{minipage}[b]{0.30\columnwidth}\raggedright
Tools and limits\strut
\end{minipage}\tabularnewline
\midrule
\endhead
\begin{minipage}[t]{0.30\columnwidth}\raggedright
Independent C1, C2-r2, C3, and B1-r2 results\strut
\end{minipage} & \begin{minipage}[t]{0.30\columnwidth}\raggedright
\href{https://github.com/MoyaMryia/riscv-accl/blob/main/spacemit/reports/2026-10-09-local-candidate-result-rejudgment.md}{Candidate
review},
\href{https://github.com/MoyaMryia/riscv-accl/blob/main/spacemit/reports/2026-10-11-project-progress.md}{progress
snapshot},
\href{https://github.com/MoyaMryia/riscv-accl/blob/main/spacemit/reports/2026-10-11-original-current-comparison.csv}{phase
comparison CSV}\strut
\end{minipage} & \begin{minipage}[t]{0.30\columnwidth}\raggedright
\href{https://github.com/MoyaMryia/riscv-accl/blob/main/spacemit/reports/2026-10-11-original-current-comparison.md}{Version
comparison}; C0 already has RVV attention and route 3; no cumulative
speedup\strut
\end{minipage}\tabularnewline
\begin{minipage}[t]{0.30\columnwidth}\raggedright
Combined-r1 fails one decode gate\strut
\end{minipage} & \begin{minipage}[t]{0.30\columnwidth}\raggedright
\href{https://github.com/MoyaMryia/riscv-accl/blob/main/spacemit/reports/2026-10-12-combined-stop-audit.md}{Stop
audit},
\href{https://github.com/MoyaMryia/riscv-accl/blob/main/spacemit/reports/raw/k1-combined-stop-audit-20261012/input/summary.json}{saved
original summary},
\href{https://github.com/MoyaMryia/riscv-accl/blob/main/spacemit/reports/raw/k1-combined-stop-audit-20261012/local-audit.json}{independent
audit}\strut
\end{minipage} & \begin{minipage}[t]{0.30\columnwidth}\raggedright
\href{https://github.com/MoyaMryia/riscv-accl/blob/main/spacemit/bench/local-candidates/audit-combined-stop.py}{Audit
tool}; 96 first-round requests and 48 comparisons; no second round or
complete-answer stage\strut
\end{minipage}\tabularnewline
\begin{minipage}[t]{0.30\columnwidth}\raggedright
Dispatch revision remains local\strut
\end{minipage} & \begin{minipage}[t]{0.30\columnwidth}\raggedright
\href{https://github.com/MoyaMryia/riscv-accl/blob/main/spacemit/experiments/2026-10-12-combined-dispatch-r2-local-protocol.md}{Protocol},
\href{https://github.com/MoyaMryia/riscv-accl/blob/main/spacemit/reports/raw/k1-combined-stop-audit-20261012/dispatch-manifest.json}{manifest},
\href{https://github.com/MoyaMryia/riscv-accl/blob/main/spacemit/reports/raw/k1-combined-stop-audit-20261012/tests.log}{local
checks}\strut
\end{minipage} & \begin{minipage}[t]{0.30\columnwidth}\raggedright
Full numerical, target assembly, operator, and model checks remain
pending\strut
\end{minipage}\tabularnewline
\begin{minipage}[t]{0.30\columnwidth}\raggedright
V7 wide intervals and incorrect point directions\strut
\end{minipage} & \begin{minipage}[t]{0.30\columnwidth}\raggedright
\href{https://github.com/MoyaMryia/riscv-accl/blob/main/spacemit/reports/2026-10-12-k1-model-v7.md}{V7
report},
\href{https://github.com/MoyaMryia/riscv-accl/blob/main/spacemit/reports/raw/k1-model-v7-20261012/evaluation.json}{evaluation},
\href{https://github.com/MoyaMryia/riscv-accl/blob/main/spacemit/reports/raw/k1-model-v7-20261012/calibration.json}{calibration},
\href{https://github.com/MoyaMryia/riscv-accl/blob/main/spacemit/reports/raw/k1-model-v7-20261012/cross-host-check.json}{cross-host
check}\strut
\end{minipage} & \begin{minipage}[t]{0.30\columnwidth}\raggedright
\href{https://github.com/MoyaMryia/riscv-accl/blob/main/spacemit/bench/qemu-screen/predictor/README.md}{Evaluator
guide}; 9/9 group coverage does not establish reliable point
predictions\strut
\end{minipage}\tabularnewline
\begin{minipage}[t]{0.30\columnwidth}\raggedright
First native operator training batch\strut
\end{minipage} & \begin{minipage}[t]{0.30\columnwidth}\raggedright
\href{https://github.com/MoyaMryia/riscv-accl/blob/main/spacemit/reports/2026-10-12-k1-operator-data.md}{Operator
report},
\href{https://github.com/MoyaMryia/riscv-accl/blob/main/spacemit/reports/raw/k1-operator-data-20261012-r1/plan.json}{plan},
\href{https://github.com/MoyaMryia/riscv-accl/blob/main/spacemit/reports/raw/k1-operator-data-20261012-r1/observations.jsonl}{observations},
\href{https://github.com/MoyaMryia/riscv-accl/blob/main/spacemit/reports/raw/k1-operator-data-20261012-r1/summary.json}{summary}\strut
\end{minipage} & \begin{minipage}[t]{0.30\columnwidth}\raggedright
\href{https://github.com/MoyaMryia/riscv-accl/blob/main/spacemit/bench/operator-data/README.md}{Operator
tools},
\href{https://github.com/MoyaMryia/riscv-accl/blob/main/spacemit/reports/raw/k1-operator-data-20261012-r1/numeric-receipts.json}{numerical
receipts},
\href{https://github.com/MoyaMryia/riscv-accl/blob/main/spacemit/reports/raw/k1-operator-data-20261012-r1/transfer-receipt.json}{transfer
receipt}; training only\strut
\end{minipage}\tabularnewline
\begin{minipage}[t]{0.30\columnwidth}\raggedright
Access failure and unknown second-batch state\strut
\end{minipage} & \begin{minipage}[t]{0.30\columnwidth}\raggedright
\href{https://github.com/MoyaMryia/riscv-accl/blob/main/spacemit/reports/2026-10-12-k1-connection-interruption.md}{Connection
report}\strut
\end{minipage} & \begin{minipage}[t]{0.30\columnwidth}\raggedright
\href{https://github.com/MoyaMryia/riscv-accl/blob/main/spacemit/bench/operator-data/remote.py}{Bounded
client}; no automatic retry or new board task\strut
\end{minipage}\tabularnewline
\bottomrule
\end{longtable}

The October 11 progress report and comparison CSV retain their original
snapshots. The October 12 stop audit supplies the later combination
decision. A saved summary is not a new independent comparison of the 536
formal state dumps. The report retains that evidence gap.

\Needspace{250pt}

\hypertarget{portable-evaluator-replay}{%
\subsection{7. Portable evaluator
replay}\label{portable-evaluator-replay}}

Run these commands from the repository root. The output directory must
not exist.

\begin{Shaded}
\begin{Highlighting}[]
\ExtensionTok{python3}\NormalTok{ spacemit/bench/qemu{-}screen/predictor/model{-}bundle.py }\KeywordTok{\textbackslash{}}
  \ExtensionTok{{-}{-}include{-}ime{-}audit}\NormalTok{ {-}{-}output lab/k1{-}portable{-}v7}
\end{Highlighting}
\end{Shaded}

Check the generated archive against its external SHA-256 receipt.
Extract the archive. Run this command from its \nolinkurl{k1-models-v7/}
directory.

\begin{Shaded}
\begin{Highlighting}[]
\ExtensionTok{python3}\NormalTok{ run{-}examples.py {-}{-}output ../k1{-}replay}
\end{Highlighting}
\end{Shaded}

The replay checks fixed evidence and frozen model identities. It does
not collect new native measurements or qualify a new algorithm. Complete
archives and binary dumps remain outside Git. Full archive checks
require those retained files. The portable examples include the data
needed for their own replay.
\end{document}
